% concatenated from LyuZhouOscillatoryJacobitype.tex
%------------------------------------------------------------------------------
% Beginning of journal.tex
%------------------------------------------------------------------------------
%
% AMS-LaTeX version 2 sample file for journals, based on amsart.cls.
%
%        ***     DO NOT USE THIS FILE AS A STARTER.      ***
%        ***  USE THE JOURNAL-SPECIFIC *.TEMPLATE FILE.  ***
%
% Replace amsart by the documentclass for the target journal, e.g., tran-l.
%
\documentclass[11pt,a4paper]{article}

\usepackage{makecell}
\usepackage{authblk}
\usepackage{amsmath,amssymb,amsthm}
\allowdisplaybreaks
\usepackage[margin=1in]{geometry}
\usepackage{color}
\usepackage{float}
\usepackage{enumerate}
\usepackage{listings}
\usepackage{xcolor}
\usepackage[title]{appendix}
\lstset{
    language=Mathematica,
    basicstyle=\ttfamily\small,
    commentstyle=\color{gray!80!black},
    breaklines=true,
    breakatwhitespace=true,
    prebreak=\space\space,
    postbreak=\space,
    breakindent=2em,
    columns=fullflexible,
    frame=single,
    framerule=0.4pt,
    framesep=2pt,
    xleftmargin=0.5em,
    xrightmargin=0.5em,
    showstringspaces=false,
    tabsize=4,
    basewidth={0.5em,0.45em},
}

\newtheorem{theorem}{Theorem}[section]
\newtheorem{lemma}[theorem]{Lemma}
\newtheorem{proposition}[theorem]{Proposition}
\newtheorem{corollary}[theorem]{Corollay}
\newtheorem{definition}[theorem]{Definition}
\newtheorem{example}[theorem]{Example}
\newtheorem{remark}[theorem]{Remark}
\newtheorem{conjecture}[theorem]{Conjecture}

\numberwithin{equation}{section}



\begin{document}

\title{Recurrence Coefficients of the Orthogonal Polynomials for Oscillatory Jacobi-type Weight Functions}

\author[1,*]{Shulin Lyu}
\author[1,$\dag$]{Xun Zhou}  % *���ͨѶ����
\affil[1]{School of Mathematics and Statistics, Qilu University of Technology (Shandong Academy of Sciences), Jinan 250353, China}
\affil[*]{\texttt{lvshulin1989@163.com}, Corresponding author}
\affil[$\dag$]{\texttt{zhouxun200212@163.com}}
\date{\today}
\maketitle

\begin{abstract}
\sloppy
We study two classes of oscillatory Jacobi-type weight functions $x^c(1-x^2)^{\lambda-1/2}\exp(i\zeta x)$, $x\in[-1,1], \lambda>-1/2, c\in\{0,1\}$. Other restrictions are imposed on $\lambda$ and $\zeta$ to guarantee the existence of the associated orthogonal polynomials. By using the ladder operators established in the recent literature for monic orthogonal polynomials associated with Jacobi-type weight functions and three compatibility conditions, we derive two coupled difference equations satisfied by the three-term recurrence coefficients. Compared with the existing results, these equations are structurally simpler and of lower order. Once the initial values are determined, the recurrence coefficients can be computed at any stage through the difference equations. The obtained expressions enable us to conjecture the symbolic forms for the recurrence coefficients.
\end{abstract}




\section{Introduction}
In this paper, we consider two classes of oscillatory Jacobi-type weight functions. The first class is the oscillatory Gegenbauer weight, i.e.
\begin{align}\label{w1}
w_1(x)=(1-x^2)^{\lambda-\frac{1}{2}}\exp(i\zeta x), \qquad x\in[-1,1],
\end{align}
where $\lambda$ is a rational number, $\lambda>-\frac{1}{2}$ and $\zeta\in\mathbb{R}\setminus\{0\}$ is a zero of the Bessel function of the first kind $J_{\lambda-1}$. Here $i$ is the imaginary unit, i.e. $i^2=-1$, and $J_{\nu}$ is defined by \cite[p. 102]{Lebedev}:
\[J_{\nu}=\sum_{k=0}^{\infty}\frac{(-1)^k(z/2)^{\nu+2k}}{k!\Gamma(\nu+k+1)}.\]
The second class is
\begin{align}\label{w2}
w_2(x) = x(1-x^{2})^{\gamma-\frac{1}{2}} \exp(\mathit{i}\eta x) ,	 \quad x\in[-1,1] ,\quad \gamma >-\frac{1}{2}.
\end{align}
The real parameters $\gamma$ and $\eta$ are chosen such that there exists a unique sequence of monic orthogonal polynomials associated with $w_2(x)$.


It was demonstrated in \cite[Theorem 2.2]{MCM} that there exists a sequence of monic polynomials orthogonal with respect to $w_1(x)$, which extends the results of \cite{MCS} where $\lambda$ is assumed to be a positive rational number. In addition, the symbolic forms of the three-term recurrence coefficients for the orthogonal polynomials were conjectured. Moreover, by using the differential relations for the orthogonal polynomials, difference equations were deduced for the recurrence coefficients. The existence for the weight function $(1-x)^{\alpha-1/2}(1+x)^{\beta-1/2}\exp(i\zeta x),x\in[-1,1]$ was established in \cite{SC}, where $\alpha,\beta>-1/2$ are rational numbers, $\beta=\alpha\pm \ell$ with $\ell$ being a positive integer, and $\zeta$ is a positive zero of $J_{\alpha-1}$. Note that the special case $\alpha=\beta=\lambda$, i.e. $\ell=0$, which reduces to \eqref{w1}, was excluded. Difference equations satisfied by the recurrence coefficients, though not presented explicitly, are obtainable from the given results in \cite{SC} by one further step (see Theorems 3.3-3.5 therein). From these equations and given initial values, the recurrence coefficients can be computed, which enables the authors of \cite{SC} to construct the effective quadrature rules of Gaussian type.

For the weight function given by \eqref{w2}, the existence of monic orthogonal polynomials was proved when $\gamma=0$ and $\eta$ is a positive zero of the Bessel function $J_0$ (see \cite[Theorem 4]{MC05-1}), and when $\gamma=1/2$ and $\eta=m\pi$ with $m$ denoting a nonzero integer (see \cite[Theorem 2.3]{MC05-2}). For the former case (i.e. $w_2(x)=x(1-x^2)^{-1/2}\exp(i\zeta x)$), as the degree of the monic orthogonal polynomial tends to $\infty$, the asymptotic formulas for the recurrence coefficients were obtained, based on the asymptotic properties for the orthogonal polynomials associated with the weight function $(1-x^2)^{-1/2}\exp(i\zeta x)$. For the recurrence coefficients of the latter case (i.e. $w_2(x)=x\exp(i\zeta x)$), symbolic forms were conjectured with the aid of a software package and difference equations were deduced by using differential identities satisfied by the orthogonal polynomials. Gaussian quadrature rules were also constructed and applied to the numerical evaluation of highly oscillatory integrals.

Orthogonal polynomials associated with weight functions containing an oscillatory factor and the computation of oscillatory integrals by using quadrature rules were investigated in several studies. In \cite{MSM}, the weight function $x^{\alpha}{\rm e}^{-x}\exp(i\zeta x), x\in[0,+\infty),\alpha>-1,\zeta>0$ was considered. The existence of the associated monic orthogonal polynomials was proved and the corresponding quadrature rules of Gaussian type was constructed. The highly oscillatory integrals with respect to $[\log(x-a)]^{c_1}[\log(b-x)]^{c_2}e^{i\omega x}(x-a)^{-\alpha}(b-x)^{-\beta}, x\in[a,b]$, $\omega\gg1, c_1,c_2\in\{0,1\}, \alpha,\beta<1$, was evaluated numerically in \cite{KHM} by using the three-term recurrence coefficients of the orthogonal polynomials associated with Gautschi's weight $t^{s}{\rm e}^{-t}(t-1-\log t), t\in(0,\infty),s>-1$ and the corresponding Gaussian-type quadrature rules. The oscillatory integrals with kernel $x^{\alpha}(b-x)^{\beta}\exp(i\omega x^{\gamma}), x\in[0,b], |\omega|\gg1, -1<\alpha,\beta\leq0, \gamma$ being a positive integer, was computed in \cite{CYC} via an $s$-step asymptotic rule and asymptotic-Filon-type quadrature rule. See \cite{KX} for analysis on four kinds of highly oscillatory Fourier-type integrals.

Denote $w_1(x)$ and $w_2(x)$ by $w(x), x\in[-1,1]$, and the associated monic orthogonal polynomials by $\{P_n(x)\}$. Write
\begin{align}\label{mp}
P_n(x):=x^n+\textbf{p}(n)x^{n-1}+\dots+P_n(0),\quad\qquad n=1,2,\cdots,
\end{align}
with $P_0(x):=1$ and $\textbf{p}(0):=0$, and suppose
\begin{align}\label{dp}	
\int_{-1}^1P_m(x)P_n(x)w(x)dx=h_n\delta_{mn},\quad\qquad m,n=0,1,\cdots,
\end{align}
where $\delta_{mn}$ is the Kronecker delta function, i.e. $\delta_{mn}=1$ for $m=n$ and $0$ otherwise. According to \eqref{mp}-\eqref{dp}, ones gets the three-term recurrence relation
\begin{align}\label{3r}
x\,P_n(x)=P_{n+1}(x)+i\alpha_nP_n(x)+\beta_nP_{n-1}(x),\qquad\qquad n\geq0,
\end{align}
with $P_{-1}(x):=0$. In much of the literature, $\beta_0$ is set to 0, but in this paper we keep it arbitrary. The recurrence coefficients  are given by
\begin{align}
i\alpha_n=&\frac{1}{h_n}\int_{-1}^1xP_n^2(x)w(x)dx\label{al-1}\\
=&\textbf{p}(n)-\textbf{p}(n+1),\qquad\qquad n\geq0,\label{al}\\
\beta_n=&\frac{1}{h_{n-1}}\int_{-1}^1xP_n(x)P_{n-1}(x)w(x)dx\label{bt-1}\\
=&\frac{h_n}{h_{n-1}},\qquad\qquad n\geq1.\label{bt}
\end{align}
Refer to \cite[p. 22]{Ismail} and \cite{ML25}. It follows from \eqref{al} that
\begin{align}
\mathit{i}\sum\limits_{j=0}^{n-1} \alpha_{j}=-\textbf{p}(n).\label{(1.7)}
\end{align}

In \cite{LyuLyu25}, the ladder operators for the monic orthogonal polynomials associated with the Jacobi-type weight function
$(1-x)^{\alpha}(1+x)^{\beta}w_0(x),~ x\in[-1,1],~\alpha,\beta>-1$,
were established. Here $w_0(x)>0$ is continuously differentiable and bounded at $\pm1$. It should be noted that the condition $w_0(x)>0$ is imposed to guarantee the existence of the orthogonal polynomials, and it was not used in the derivation of the ladder operators and compatibility conditions. Therefore, Theorem 4.1 of \cite{LyuLyu25} also applies to the weight function $w(x)$ given by \eqref{w1} and \eqref{w2}. We restate it below.
\begin{proposition}
The monic orthogonal polynomials $\{P_n\}$ defined by \eqref{mp}-\eqref{dp} satisfy the following ladder operators:
\begin{align}
P_n'(z)=&-B_n(z)P_n(z)+\beta_nA_n(z)P_{n-1}(z),\label{loJ}\\
P_{n-1}'(z)=&\left(B_n(z)+v'(z)\right)P_{n-1}(z)-A_{n-1}(z)P_n(z),\label{roJ}
\end{align}
for $n\geq0$, where $v(x):=-\ln w(x)$. The quantities $A_n(z)$ and $B_n(z)$ are defined by
\begin{align}\label{An}
A_n(z):=\frac{1}{1-z^2}\left[  \frac{1}{h_n}\int_{-1}^{1}\frac{(1-z^2)v'(z)-(1-x^2)v'(x)}{z-x}P_n^2(x)w(x)dx+2n+1 \right],
\end{align}
for $n\geq0$ with $A_{-1}(z):=0$, and
\begin{align}\label{Bn}		 B_n(z):=&\frac{1}{1-z^2}\cdot\frac{1}{h_{n-1}}\int_{-1}^{1}\frac{(1-z^2)v'(z)-(1-x^2)v'(x)}{z-x}P_n(x)P_{n-1}(x)w(x)dx+\frac{nz-\textbf{p}(n)}{1-z^2}, \end{align}
for $n\geq1$ with $B_0(z):=0$. Furthermore, $A_n(z)$ and $B_n(z)$ satisfy the following three compatibility conditions
 \begin{align}
B_{n+1}(z)+B_n(z)&=(z-i\alpha_n)A_n(z)-v'(z)\tag{$S_1$},\\
1+(z-i\alpha_n)(B_{n+1}(z)-B_n(z))&=\beta_{n+1}A_{n+1}(z)-\beta_nA_{n-1}(z)\tag{$S_2$},			 \end{align}
for $n\geq0$, and
\begin{align}
B^2_n(z)+v'(z)B_n(z)+\sum_{j=0}^{n-1}A_j(z)&=\beta_nA_n(z)A_{n-1}(z)\tag{$S^{'}_2$},
\end{align}
for $n\geq1$.
\end{proposition}

The ladder operators and $\{(S_1), (S_2), (S_2')\}$ were widely used in literature to study real nonnegative Jacobi-type weight functions which have finite moments of all orders and contain one or more variables, including the associated monic orthogonal polynomials, the partition functions of the unitary ensembles, the gap probabilities, etc. See, for instance, \cite{BCH, CI, MC17, MC20, MC21}. It should be noted that, in previous studies, $\alpha$ and $\beta$ are usually assumed to be positive, and the expressions of $A_n$ and $B_n$ are different from \eqref{An}-\eqref{Bn}. The treatment for the Laguerre-type weight functions $x^{\lambda}w_0(x), x\in[0,+\infty)$, is analogous. See, e.g., \cite{LyuLyu25, MF25} for the derivation of the ladder operators with $\lambda>-1$ and \cite{CI10, CM12, ML25} for the applications of the ladder operators to problems with $\lambda>0$. Using the difference and Toda-type equations established for the recurrence coefficients via the ladder operators, researchers deduced large-$n$ and large-time asymptotics for the recurrence coefficients and related quantities. We refer to recent work, e.g., \cite{CJ21,MF25, MW26-1, MW26-2}.


In this paper, we apply the ladder operators \eqref{loJ}-\eqref{roJ} and $\{(S_1), (S_2), (S_2')\}$ to the two types of Jacobi-type weight functions given by \eqref{w1} and \eqref{w2}, and deduce difference equations for the recurrence coefficients $\alpha_n$ and $\beta_n$.
Our derivation process is outlined below. We first calculate $A_n$ and $B_n$ using \eqref{An} and \eqref{Bn}, which are expressed in terms of $\{\alpha_n,\beta_n\}$ and auxiliary quantities. Then we substitute the obtained expressions into $\{(S_1), (S_2')\}$ for $w_1(x)$ and into $\{(S_1),(S_2)\}$ for $w_2(x)$, and obtain a series of difference equations. Eliminating the auxiliary quantities from these equations, we finally arrive at a system of two coupled difference equations for $\alpha_n$ and $\beta_n$. The complete derivations for our weight function $w_1(x)$ and $w_2(x)$ are given in Sections 2 and 3, respectively. Our conclusions are presented in Section 4.
\section{Oscillatory-Gegenbauer Weight and Difference Equations}
For the weight function given by \eqref{w1}, we have
\[v(x)=-\ln w_1(x)=-i\zeta x-\left(\lambda-\frac{1}{2}\right)\ln (1-x^2).\]
It follows that
\[\frac{(1-z^2)v'(z)-(1-x^2)v'(x)}{z-x}=i\zeta x+i\zeta z+2\lambda-1.\]
Plugging it into \eqref{An}-\eqref{Bn}, in view of the orthogonality relation \eqref{dp} and \eqref{al}-\eqref{bt}, we come to the expressions for $A_n$ and $B_n$.
\begin{lemma}
$A_n(z)$ and $B_n(z)$ are given by
\begin{align}
A_n(z)=&\frac{i\zeta z-\zeta\alpha_n+2(n+\lambda)}{1-z^2},\qquad\qquad n\geq0,\label{An-J}\\
B_n(z)=&\frac{nz+i\zeta\beta_n-\textbf{p}(n)}{1-z^2},\qquad\qquad n\geq1.\label{Bn-J}
\end{align}
\end{lemma}

Substituting \eqref{An-J}-\eqref{Bn-J}, together with the initial values $A_{-1}(z)=0$ and $B_0(z)=0$, into $(S_1)$ and $(S_2')$, we obtain two difference equations for $\alpha_n$ and $\beta_n$.
\begin{theorem} The recurrence coefficients $\alpha_n$ and $\beta_n$ satisfy the following two coupled difference equations
\begin{align}
\alpha_{n+1}=&\frac{3}{2\zeta}\left(2n+1+2\lambda\right)-\alpha_{n}+\frac{2n+1+2\lambda}{2\zeta\beta_{n+1}}\left[\beta_{n}-1+\alpha_{n}\left(\frac{2n-1+2\lambda}{\zeta}-\alpha_{n}\right)\right],\label{diff2}\\
\beta_{n+2}=&\beta_n+\alpha_{n+1}^2-\alpha_n^2+\frac{1}{\zeta}\big[\alpha_n(2n-1+2\lambda)-\alpha_{n+1}(2n+3+2\lambda)\big],\label{diff1}
\end{align}
for $n\geq1$, with the initial conditions given by
\begin{subequations}\label{al0bt1}
\begin{align}
\alpha_0=&\frac{2\lambda}{\zeta},&\alpha_1=&\frac{2(\lambda+1)}{\zeta}-\frac{\zeta(2\lambda+1)}{\zeta^2-2\lambda},\\
\beta_{1}=&1-\frac{2\lambda}{\zeta^2}, &\beta_2=&-\frac{2(2\lambda+1)\left(\zeta ^4-\zeta ^2 \lambda  (2 \lambda +5)+4 \lambda ^2\right)}{\zeta^2(\zeta^2-2\lambda)^2}.
\end{align}
\end{subequations}
Note that $\beta_0$ can be arbitrary. It was specified in \cite{MCM} as $\sqrt{\pi}~\Gamma\left(\lambda+\frac{1}{2}\right)\left(\frac{2}{\zeta}\right)^{\lambda}J_{\lambda}(\zeta)$, whereas in much of the literature it is taken as zero. In this paper, we leave it arbitrary.
\end{theorem}
\begin{proof}
Inserting \eqref{An-J}-\eqref{Bn-J} into $(S_1)$ gives us
\begin{align}\label{S1}
i\zeta\left(\beta_{n+1}+\beta_n\right)-\textbf{p}(n+1)-\textbf{p}(n)=i\alpha_n\left(\zeta\alpha_n-2(n+\lambda)\right)+i\zeta.
\end{align}
According to \eqref{al}, we replace $\textbf{p}(n+1)$ by $\textbf{p}(n)-i\alpha_n$ in the above identity and get
\begin{align}\label{S1-1}
\textbf{p}(n)=\frac{1}{2}\left[i\zeta(\beta_{n+1}+\beta_n-1)+i\alpha_n(-\zeta\alpha_n+2n+1+2\lambda)\right].
\end{align}
Replacing $n$ by $n+1$ in this equality and subtracting the resulting equation from \eqref{S1-1}, in light of the fact that $\textbf{p}(n)-\textbf{p}(n+1)=i\alpha_n$ given by \eqref{al}, we are led to \eqref{diff1}.

To continue, we plug \eqref{An-J}-\eqref{Bn-J} into $(S_2')$, multiply both sides of the obtained identity by $\left(1-z^2\right)^2$ and find
\begin{equation*}
\begin{aligned}
&(2n-1+2\lambda)\left(i\zeta\beta_n-\textbf{p}(n)\right)\cdot z+(i\zeta\beta_n-\textbf{p}(n))^2+n(n-1+2\lambda)\\
&=i\zeta\beta_n\big(-\zeta(\alpha_{n-1}+\alpha_n)+2(2n-1+2\lambda)\big)\cdot z\\
&\quad+\beta_n\big(-\zeta\alpha_n+2(n+\lambda)\big)\big(-\zeta\alpha_{n-1}+2(n-1+\lambda)\big)-\zeta^2\beta_n.
\end{aligned}
\end{equation*}
Comparing the coefficient of $z$ on both sides yields
\begin{align}
(2n-1+2\lambda)\left(i\zeta\beta_n+\textbf{p}(n)\right)=&i\zeta^2\beta_n\big(\alpha_{n-1}+\alpha_n\big).\label{S2'-1}
%(i\zeta\beta_n-\textbf{p}(n))^2+n(n-1+2\lambda)=&\beta_n\big(-\zeta\alpha_n+2(n+\lambda)\big)\big(-\zeta\alpha_{n-1}+2(n-1+\lambda)\big)-\zeta^2\beta_n.\label{S2'-2}
\end{align}
Substituting \eqref{S1-1} into \eqref{S2'-1}, with $n$ replaced by $n+1$ in the resulting equation, we come to
\begin{align*}
\beta_{n+2}=\frac{2\zeta\beta_{n+1}\left(\alpha_n+\alpha_{n+1}\right)}{2n+1+2\lambda}+\alpha_{n+1}\left(\alpha_{n+1}-\frac{2n+3+2\lambda}{\zeta}\right)-3\beta_{n+1}+1.
\end{align*}
Combining it with \eqref{diff1} leads us to \eqref{diff2}.

The detailed computations of $\{\alpha_0,\alpha_1,\beta_1,\beta_2\}$ are presented in Appendix \ref{J0D-int-1}. Alternatively, their expressions can be obtained using the Mathematica code provided in Appendix \ref{IV1-M}.
\end{proof}

\begin{remark}
Two coupled difference equations were also deduced for $\alpha_n$ and $\beta_n$ in \cite[Theorem 3.5]{MCM}. One is the same as \eqref{diff1}, and the other is of second order in both $\alpha_n$ and $\beta_n$, which is more complicated than the first-order difference equation \eqref{diff2}. Consequently, given $\{\alpha_0,\alpha_1,\beta_1,\beta_2\}$ as in \eqref{al0bt1}, $\alpha_n$ and $\beta_n$ for any $n$ can be computed using \eqref{diff2}-\eqref{diff1}, which are independent of the value of $\beta_0$. However, to achieve this goal by using the difference equations in \cite{MCM}, the initial values $\{\alpha_1,\alpha_2,\beta_1,\beta_2\}$ are needed.
\end{remark}

\begin{corollary}
Since the initial values given by \eqref{al0bt1} and the coefficients that appear in the difference equations \eqref{diff2}-\eqref{diff1} are all real, we conclude that $\{\alpha_n,n\geq0\}$ and $\{\beta_n, n\geq1\}$ are real, which coincides with Lemma 3.1 of \cite{MCS} where the parameter $\lambda$ in the weight function was assumed to be positive.
\end{corollary}

Iterating the difference equations \eqref{diff2} and \eqref{diff1} for $1\leq n\leq7$ yields $\alpha_i$ for $2\leq i\leq8$ and $\beta_i$ for $3\leq i\leq9$. The obtained results together with the initial conditions given by \eqref{al0bt1} are listed in Table \ref{Table-1}. Our expressions for $\{\alpha_i,\beta_i, 0\leq i\leq 2\}$ are identical to those reported in \cite[Table 1]{MCM}, although we were unable to trace the software package that the authors used for their computation and referred to in the cited literature.
\begin{table}[H]
\centering
\renewcommand{\arraystretch}{1.2}
\setlength{\tabcolsep}{3pt}
\caption{Recurrence coefficients $\{\alpha_n\}_{n=0}^8$ and $\{\beta_n\}_{n=0}^9$ with arbitrary $\lambda$ and  $\beta_0$}\label{Table-1}
\begin{tabular}{|c|c|c|}
\hline
$n$ & $\alpha_n$ & $\beta_n$\\
\hline
0 & $\frac{2\lambda}{\zeta}$ & $\beta_{0}$ \\
\hline
1 & $\frac{2(\lambda+1)}{\zeta}-\frac{\zeta(2\lambda+1)}{\zeta^2-2\lambda}$ & $\frac{\zeta^2-2\lambda}{\zeta^2}$ \\
\hline
2 & $\frac{2(\lambda+2)}{\zeta}+\zeta\left(\frac{2\lambda+1}{\zeta^2-2\lambda}-\frac{(2\lambda+3) (\zeta^2-4\lambda^2-4\lambda)}{2X_{2}}\right)$ & $-\frac{2(2\lambda+1)X_{2}}{\zeta^2(\zeta^2-2\lambda)^2}$ \\
\hline
3 & $\frac{2(\lambda+3)}{\zeta}+(2\lambda+3)\zeta\left(\frac{ \zeta^2-4\lambda^2-4\lambda}{2X_{2}}-\frac{2X_{3}}{X_{4}}\right)$ & $\frac{(\zeta^2-2\lambda)X_{4}}{\zeta^2X_{2}^2}$  \\
\hline
4 &$\frac{2(\lambda+4)}{\zeta}+\zeta\left(\frac{2(2\lambda+3)X_{3}}{X_{4}}-\frac{(2\lambda+5)X_{5}}{X_{6}}\right)$&$-\frac{4(2\lambda+3)X_{2}X_{6}}{\zeta^2X_{4}^2}$ \\
\hline
5 &$\frac{2(\lambda+5)}{\zeta}+(2\lambda+5)\zeta\left(\frac{X_{5}}{X_{6}}-\frac{3X_{8}}{X_{9}}\right)$&$\frac{X_{4}X_{9}}{\zeta^2X_{6}^2}$ \\
\hline
6 &$\frac{2(\lambda+6)}{\zeta}+3\zeta\left(\frac{(2\lambda+5)X_{8}}{X_{9}}-\frac{(2\lambda+7)X_{11}}{2X_{12}}\right)
$& $-\frac{6(2\lambda+5)X_{6}X_{12}}{\zeta^2X_{9}^{2}}$\\
\hline
7&$\frac{2(\lambda+7)}{\zeta}+(2\lambda+7)\zeta\left(\frac{3X_{11}}{2X_{12}}-\frac{4X_{15}}{X_{16}}\right)$&$\frac{X_{9}X_{16}}{\zeta^2X_{12}^{2}}$\\
\hline
8&$\frac{2(\lambda+8)}{\zeta}+2\zeta\left[\frac{2(2\lambda+7)X_{15}}{X_{16}}-\frac{(2\lambda+9)X_{19}}{X_{20}}\right]$&$-\frac{8(2\lambda+7)X_{12}X_{20}}{\zeta^2X_{16}^{2}}$\\\hline
9&-&$\frac{X_{16}X_{25}}{\zeta^2X_{20}^{2}}$\\\hline
\end{tabular}
\vspace{0.4em}
\par {\small Note: $X_i=X_i(\zeta^2)$, with $X_i(\cdot)$ denoting a monic polynomial of degree $i$. Hence,\\ $X_i(\zeta^2)$ is a monic polynomial of degree $2i$ in $\zeta$, containing only even powers of $\zeta$.\\ The expressions of $X_i(\zeta^2)$ for $i=2,3,4,5,6,8,9$ are given in Appendix \ref{symbols},\\
while those for $i=11,12,15,16,19,20,25$ are not presented for brevity.}
\end{table}

According to Table \ref{Table-1}, we propose the following conjecture.
\begin{conjecture}
The recurrence coefficients $\alpha_n$ with $n\geq4$ have the following symbolic forms:
\begin{align*}
\alpha_{2k}=&\frac{2(\lambda+2k)}{\zeta}+k\zeta\left((2\lambda+2k-1)\frac{X_{k^2-1}(\zeta^2)}{X_{k^2}(\zeta^2)}-(2\lambda+2k+1)\frac{X_{k(k+1)-1}(\zeta^2)}{2X_{k(k+1)}(\zeta^2)}\right),\\ \alpha_{2k+1}=&\frac{2(\lambda+2k+1)}{\zeta}+(2\lambda+2k+1)\zeta\left(k\frac{X_{k(k+1)-1}(\zeta^2)}{2X_{k(k+1)}(\zeta^2)}-(k+1)\frac{X_{k(k+2)}(\zeta^2)}{X_{(k+1)^2}(\zeta^2)}\right),
\end{align*}
for $k\geq2$. Regarding $\beta_n$ with $n\geq 1$, we have
\begin{align*}
\beta_{2k} =& -2k(2\lambda+2k-1)\frac{X_{k(k-1)}(\zeta^2)\cdot X_{k(k+1)}(\zeta^2)}{\zeta^2 \left(X_{k^2}(\zeta^2)\right)^2},& \beta_{2k+1} =& \frac{X_{k^2}(\zeta^2)\cdot X_{(k+1)^2}(\zeta^2)}{\zeta^2 \left(X_{k(k+1)}(\zeta^2)\right)^2},
\end{align*}
which hold for $k\geq0$, except for $k=0$ in $\beta_{2k}$. Here $X_j(\cdot)$ is a monic polynomial of degree $j$ for $j\geq0$, with $X_0(\zeta^2):=1$ and $X_1(\zeta^2):=\zeta^2-2\lambda$. Note that $X_i(\zeta^2)$ for $i=2,3,4,5,6,8,9$ are listed in Appendix \ref{symbols}.
\end{conjecture}
\begin{remark}
The symbolic forms of the recurrence coefficients are also conjectured in \cite{MCM} (below Table 1 therein), according to the calculations of the Hankel determinant $H_n$. We rewrite them as follows:
\begin{align*}
\alpha_{2k} =& \frac{V_{k(2k+1)}(\zeta^2)}{\zeta S_{k^2}(\zeta^2)\cdot S_{k(k+1)}(\zeta^2)}, & \beta_{2k} =& -\frac{S_{k(k-1)}(\zeta^2)\cdot S_{k(k+1)}(\zeta^2)}{\zeta^2 \left(S_{k^2}(\zeta^2)\right)^2}, \\
\alpha_{2k+1} =& \frac{V_{(k+1)(2k+1)}(\zeta^2)}{\zeta S_{k(k+1)}(\zeta^2)\cdot S_{(k+1)^2}(\zeta^2)},& \beta_{2k+1} =& -\frac{S_{k^2}(\zeta^2)\cdot S_{(k+1)^2}(\zeta^2)}{\zeta^2 \left(S_{k(k+1)}(\zeta^2)\right)^2},
\end{align*}
which hold for $k\geq0$, except for $k=0$ in $\beta_{2k}$. Here $S_j(\cdot)$ and $V_j(\cdot)$ are polynomials of degree $j$, but their leading coefficients were not specified. We find from \cite{MCM} that $S_0=H_1=1$.
The expressions of $\alpha_n$ and $\beta_n$ presented in Table \ref{Table-1} are consistent with the above conjecture, where
\begin{align*}
S_1(\zeta^2)=&-(\zeta^2-2\lambda),&S_2(\zeta^2)=&2\lambda_1X_2,\\
S_4(\zeta^2)=&4\lambda_1^2X_4,&S_6(\zeta^2)=&32 \lambda_1^3\lambda_3X_6,\\ S_9(\zeta^2)=&-256\lambda_1^4\lambda_3^2X_9,&S_{12}(\zeta^2)=&12288\lambda_1^5\lambda_3^3\lambda_5X_{12},\\
S_{16}(\zeta^2)=&589824\lambda_1^6\lambda_3^4\lambda_5^2X_{16},&S_{20}(\zeta^2)=&226492416\lambda_1^7\lambda_3^5\lambda_5^3\lambda_7X_{20},\\
S_{25}(\zeta^2)=&-86973087744\lambda_1^8\lambda_3^6\lambda_5^4\lambda_7^2X_{25},&&
\end{align*}
with $\lambda_j:=2\lambda+j$, and
\begin{align*}
V_0(\zeta^2)=&2\lambda, \qquad\qquad\qquad V_1(\zeta^2)=-(\zeta^2-4\lambda^2-4\lambda),&\\
V_3(\zeta^2)=&-(2\lambda+1)\left( (6 \lambda +7)\zeta ^6-4 \lambda (\lambda+4)(2 \lambda+3)\zeta^4+32 \lambda ^2
   (2\lambda +3)\zeta ^2-32 \lambda ^3 (\lambda +2)\right).&
\end{align*}
The explicit expressions for $V_j$ with $j=6,10,15,21,28,36$ are omitted here for brevity.
From the above expressions for $S_j$, we conjecture that
\begin{align*}
S_{k^2}(\zeta^2)=&(-1)^k 2^{k(k-1)}\prod_{j=1}^{k-1}\left(j(2\lambda+2j-1)\right)^{2(k-j)}\cdot X_{k^2}(\zeta^2),\\
S_{k(k+1)}(\zeta^2)=& 2^{k^2}\prod_{j=1}^{k}\left(j(2\lambda+2j-1)\right)^{2(k-j)+1}\cdot X_{k(k+1)}(\zeta^2),
\end{align*}
which holds for $k\geq1$, except for $k=1$ in $S_{k^2}$.
\end{remark}

When $\lambda=0$, the weight function becomes $w(x)=(1-x^2)^{-\frac{1}{2}}\exp(\mathit{i}\zeta x)$, the values of ${\alpha_{n},\beta_{n}, 0 \le n \le 13}$ are given in Table \ref{Table-1-0}.
\begin{table}[H]
	\centering
	\renewcommand{\arraystretch}{1.2}
	\setlength{\tabcolsep}{3pt}
\caption{Recurrence coefficients $\{\alpha_n,\beta_n\}_{n=0}^{13}$ for $\lambda=0$ with arbitrary $\beta_0$}\label{Table-1-0}
	\begin{tabular}{|c|c|c|}
\hline
$n$ &  $\alpha_n $&$\beta_n $  \\
\hline
0  & $0$ &$\beta_{0}$\\
\hline
1  & $\frac{1}{\zeta}$ &1\\
\hline
2  & $\frac{7}{2\zeta}$ & $-\frac{2}{\zeta^2}$\\
\hline
3  & $\frac{3^3\cdot3}{14\zeta}-\frac{30\zeta}{7(\zeta^2-21/4)}$& $\frac{\zeta^2-21/4}{\zeta^2}$ \\
\hline
4 & $\frac{2^5\cdot19}{7\cdot11\zeta}+5\zeta\left(\frac{6}{7(\zeta^2-21/4)}-\frac{7 (\zeta ^2-177/8)}{11Y_{2}}\right)$&$-\frac{12Y_{2}}{\zeta^2(\zeta^2-21/4)^2}$\\
\hline
5&$\frac{5^3\cdot7}{8\cdot 11\zeta}+35\zeta\left(\frac{\zeta^2-177/8}{11Y_{2}}-\frac{3Y_{3}}{8Y_{4}}\right)$&$\frac{(\zeta^2-21/4)Y_{4}}{\zeta^2Y_{2}^2}$ \\
\hline
6&$\frac{3^4\cdot13}{8\cdot11\zeta}+21\zeta\left(\frac{5Y_{3}}{8Y_{4}}-\frac{9Y_{5}}{22Y_{6}}\right)$&$-\frac{30Y_{2}Y_{6}}{\zeta^2Y_{4}^2}$ \\
\hline
7&$\frac{7^3\cdot13}{11\cdot29\zeta}+63\zeta\left(\frac{3Y_{5}}{22Y_{6}}-\frac{12Y_{8}}{29Y_{9}}\right)$&$\frac{Y_{4}Y_{9}}{\zeta^2Y_{6}^2}$ \\
\hline
8&$\frac{4^4\cdot67}{29\cdot37\zeta}+27\zeta\left(\frac{28Y_{8}}{29Y_{9}}-\frac{22Y_{11}}{37Y_{12}}\right)$&$-\frac{56Y_{6}Y_{12}}{\zeta^2Y_{9}^{2}}$\\\hline
9&$\frac{9^3\cdot21}{23 \cdot 37\zeta}+198\zeta\left(\frac{3Y_{11}}{37Y_{12}}-\frac{5Y_{15}}{23Y_{16}}\right)$&$\frac{Y_{9}Y_{16}}{\zeta^2Y_{12}^{2}}$\\\hline
10&$\frac{5^3\cdot103}{23\cdot28\zeta}+5\zeta\left(\frac{198Y_{15}}{23Y_{16}}-\frac{143Y_{19}}{28Y_{20}}\right)$&$-\frac{90Y_{12}Y_{20}}{\zeta^2Y_{16}^{2}}$\\\hline
11&$\frac{11^3\cdot31}{28\cdot67\zeta}+715\zeta\left(\frac{Y_{19}}{28Y_{20}}-\frac{6Y_{24}}{67Y_{25}}\right)$&$\frac{Y_{16}Y_{25}}{\zeta^2Y_{20}^{2}}$\\\hline
12&$\frac{6^3\cdot588}{67\cdot79\zeta}+195\zeta\left(\frac{22Y_{24}}{67Y_{25}}-\frac{15Y_{29}}{79Y_{30}}\right)$&$-\frac{132Y_{20}Y_{30}}{\zeta^2Y_{25}^2}$\\\hline
13&$\frac{13^3\cdot43}{46\cdot79\zeta}+585\zeta\left(\frac{5Y_{29}}{79Y_{30}}-\frac{7Y_{35}}{46Y_{36}}\right)$&$\frac{Y_{25}Y_{36}}{\zeta^2Y_{30}^2}$\\\hline
\end{tabular}
\vspace{0.4em}
\par {\small Note: $Y_i=Y_i(\zeta^2)$, with $Y_i(\cdot)$ denoting a monic polynomial of degree $i$. Hence,\\ $Y_i(\zeta^2)$ is a monic polynomial of degree $2i$ in $\zeta$, containing only even powers of $\zeta$.\\ The expressions of $Y_i(\zeta^2)$ for $i=2,3,4,5,6,8,9,11,12,15,16$ are given in Appendix \ref{symbols},\\
while the others are omitted due to their length.}
\end{table}

\begin{conjecture}
The recurrence coefficients for the monic orthogonal polynomials associated with the weight function $(1-x^2)^{-1/2}\exp(i\zeta x), x\in[-1,1]$, with $\zeta\in\mathbb{R}\setminus\{0\}$ being a zero of $J_{-1}$, have the following symbolic forms:
\begin{align*}
\alpha_{2k}=&\frac{4k^3(4k^2+3)}{(k(2k-1)+1)(k(2k+1)+1)\zeta}\\
&+k(k-1)(2k+1)\zeta\left(\frac{2k-1}{k(2k-1)+1}\cdot\frac{Y_{(k-1)^2-1}(\zeta^2)}{Y_{(k-1)^2}(\zeta^2)}-\frac{2k+3}{2(k(2k+1)+1)}\cdot\frac{Y_{k(k-1)-1}(\zeta^2)}{Y_{k(k-1)}(\zeta^2)}\right), \end{align*}
for $k\geq3$,
\begin{align*}
\alpha_{4j-1}=&\frac{(4j-1)^3(2j(2j-1)+1)}{(4j^2-3j+1)(2j(4j-1)+1)\zeta}\\
&+(2j-1)(4j-1)(4j+1)\zeta\left(\frac{j-1 }{(2j-1)(4j-1)+1}\cdot\frac{Y_{2(j-1)(2j-1)-1}(\zeta^2)}{Y_{2(j-1)(2j-1)}(\zeta^2)}\right.\\
&\left.\qquad\qquad\qquad\qquad\qquad\qquad\quad-\frac{2j }{2j(4j-1)+1}\cdot\frac{Y_{(2j-1)^2-1}(\zeta^2)}{Y_{(2j-1)^2}(\zeta^2)}\right),
\end{align*}
\begin{align*}
\alpha_{4j+1}=&\frac{(4j+1)^3(2j(2j+1)+1)}{(4j^2+3j+1)(2j(4j+1)+1)\zeta}\\
&+(4j+1)(4j+3)\zeta \left(\frac{j(2j-1) }{2j(4j+1)+1}\cdot\frac{Y_{2j(2j-1)-1}(\zeta^2)}{Y_{2j(2j-1)}(\zeta^2)}-\frac{2j(2j+1) }{(2j+1)(4j+1)+1}\cdot\frac{Y_{(2j)^2-1}(\zeta^2)}{Y_{(2j)^2}(\zeta^2)}\right),
\end{align*}
for $j\geq2$, and
\begin{align*}
\beta_{2k} =& -2k(2k-1)\frac{Y_{(k-1)(k-2)}(\zeta^2)\cdot Y_{k(k-1)}(\zeta^2)}{\zeta^2 \left(Y_{(k-1)^2}(\zeta^2)\right)^2},& \beta_{2k+1} =& \frac{Y_{(k-1)^2}(\zeta^2)\cdot Y_{k^2}(\zeta^2)}{\zeta^2 \left(Y_{k(k-1)}(\zeta^2)\right)^2},
\end{align*}
for $k\geq1$ with $Y_1(\zeta^2):=\zeta^2-21/4$.
\end{conjecture}
%{\color{blue}\begin{remark}
%The calculation of $\lambda=1, \zeta=...$ agrees with one in the literature.
%\end{remark}}



\section{Oscillatory Jacobi-type Weight  and Difference Equations}
For the weight function given by \eqref{w2}, we have
\begin{align}
v(x)=-\ln w_2(x)=-\ln x-\left(\gamma-\frac{1}{2}\right)\ln (1-x^{2})-\mathit{i}\eta x,
\end{align}
so that
\begin{align}
v'(x)=-\frac{1}{x}+\frac{(2\gamma-1)x}{1-x^{2}}-\mathit{i}\eta.
\end{align}
Consequently,
\begin{align*}
\frac{(1-z^{2})v'(z)-(1-x^{2})v'(x)}{z-x}=\frac{1}{zx}+2\gamma+ \mathit{i}\eta x+\mathit{i}\eta z.
\end{align*}
Substituting it into \eqref{An} and \eqref{Bn} , we get the expressions for $ A_{n}(z) $ and $ B_{n}(z) $ .
\begin{lemma}
We have
	\begin{align}
A_{n}(z)=&\frac{\mathit{i}\eta z+2n+2\gamma+1-\eta\alpha_n}{1-z^{2}} +\frac{R_{n}(\zeta)}{z(1-z^{2})},\label{(2.8)}\\
B_{n}(z)=&\frac{nz+\mathit{i}\eta\beta_n-\textbf{p}(n)}{1-z^{2}} +\frac{r_{n}(\zeta)}{z(1-z^{2})}, \label{(2.9)}
\end{align}
where
\begin{align}
R_{n}(\zeta):=&\frac{1}{h_{n}}\int_{-1}^{1}P_{n}^2(x)(1-x^{2})^{\gamma-\frac{1}{2}}\exp(\mathit{i}\eta x)\,dx ,\qquad n\geq0,\\
r_{n}(\zeta):=&\frac{1}{h_{n-1}}\int_{-1}^{1}P_{n}(x)P_{n-1}(x)(1-x^{2})^{\gamma-\frac{1}{2}}\exp(\mathit{i}\eta x)\,dx,\qquad n\geq1.
\end{align}
\end{lemma}

Inserting \eqref{(2.8)} and \eqref{(2.9)} into {$(S_1)$}, and equating the coefficients of
$z^{-1}$, $(z-1)^{-1}$ and $(z+1)^{-1}$ on both sides, we obtain
\begin{gather}
r_{n+1}+r_{n}=1-\mathit{i}\alpha_{n}R_{n},\label{(2.12)}
\end{gather}
\begin{gather}
	\begin{split}
&-\mathit{i}\eta(\beta_{n+1}+\beta_{n})+\textbf{p}(n+1)+\textbf{p}(n)-2n-1-r_{n+1}-r_{n}\\
&\qquad\qquad=2\gamma-1-(1-\mathit{i}\alpha_{n})(2n+2\gamma+1-\eta\alpha_{n}+\mathit{i}\eta+R_{n}),\label{(2.13)}
	\end{split}
\end{gather}
\begin{gather}
	\begin{split}
&\mathit{i}\eta(\beta_{n+1}+\beta_{n})-\textbf{p}(n+1)-\textbf{p}(n)-2n-1-r_{n+1}-r_{n}\\
&\qquad\qquad=2\gamma-1-(1+\mathit{i}\alpha_{n})(2n+2\gamma+1-\eta\alpha_{n}-\mathit{i}\eta-R_{n}).\label{(2.14)}
	\end{split}
\end{gather}
From $(S_2)$, we find
\begin{gather}
-\mathit{i}\alpha_{n}(r_{n+1}-r_{n})=\beta_{n+1}R_{n+1}-\beta_{n}R_{n-1},\label{(2.16)}
\end{gather}
\begin{gather}
\begin{split}
&\left(1-\mathit{i}\alpha_{n}\right)\left[\mathit{i}\eta(\beta_{n+1}-\beta_{n})-\textbf{p}(n+1)+\textbf{p}(n)+r_{n+1}-r_{n}+1\right]\\
&\;=\beta_{n+1}\left(2n+2\gamma+3-\eta\alpha_{n+1}+\mathit{i}\eta+R_{n+1}\right)-\beta_{n}\left(2n+2\gamma-1-\eta\alpha_{n-1}+\mathit{i}\eta+R_{n-1}\right),\label{(2.17)}
\end{split}
\end{gather}
\begin{gather}
	\begin{split}
&-\left(1+\mathit{i}\alpha_{n}\right)\left[\mathit{i}\eta(\beta_{n+1}-\beta_{n})-\textbf{p}(n+1)+\textbf{p}(n)-r_{n+1}+r_{n}-1\right]\\
&\;=\beta_{n+1}\left(2n+2\gamma+3-\eta\alpha_{n+1}-\mathit{i}\eta-R_{n+1}\right)-\beta_{n}\left(2n+2\gamma-1-\eta\alpha_{n-1}-\mathit{i}\eta-R_{n-1}\right).\label{(2.18)}
\end{split}
\end{gather}

Using \eqref{(2.12)}-\eqref{(2.14)} and \eqref{(2.17)}-\eqref{(2.18)}, we express $R_n$ and $r_n$ in terms of the recurrence coefficients and $\textbf{p}(n)$.

\begin{lemma}
	The auxiliary quantities $R_{n}$ and $r_{n}$ admit the following expressions
	\begin{align}
R_{n}=&\mathit{i}\alpha_{n}(2n+2\gamma+2-\eta\alpha_{n})+\mathit{i}\eta(\beta_{n+1}+\beta_{n}-1)-2\textbf{p}(n), \label{(3.1)}
\end{align}
\begin{align}
		\begin{split}
2r_{n}=&2+2\mathit{i}\alpha_{n}\textbf{p}(n)+\alpha_{n}^{2}(2n+2\gamma+3-\eta\alpha_{n})-\eta\alpha_{n}\\
&-\beta_{n+1}(2n+2\gamma+3-\eta\alpha_{n+1}-2\eta\alpha_{n})+\beta_{n}(2n+2\gamma-1-\eta\alpha_{n-1}).\label{(3.2)}
		\end{split}
	\end{align}
\end{lemma}

\begin{proof}
Subtracting \eqref{(2.14)} from \eqref{(2.13)}, and replacing $\textbf{p}(n+1)$ in the resulting identity using \eqref{al}, we come to \eqref{(3.1)}.

Adding \eqref{(2.17)} and \eqref{(2.18)}, and using \eqref{al} to eliminate $\textbf{p}(n+1)$ in the obtained equation, we get
\begin{align}\label{s2-4}
r_{n+1}-r_{n}=\beta_{n+1}\left(2n+2\gamma+3-\eta\alpha_{n+1}-\eta\alpha_{n}\right)-\beta_{n}\left(2n+2\gamma-1-\eta\alpha_{n}-\eta\alpha_{n-1}\right)-\alpha_{n}^{2}-1.
\end{align}	
The subtraction of it from \eqref{(2.12)} where $R_n$ is replaced using \eqref{(3.1)} leads us to \eqref{(3.2)}.
\end{proof}

Using \eqref{(3.1)}-\eqref{(3.2)} to eliminate the auxiliary quantities in \eqref{(2.12)} and in the equation obtained by substituting \eqref{s2-4} into \eqref{(2.16)}, we arrive at the following result.
	\begin{proposition}
		$\textbf{p}(n)$ and the recurrence coefficients satisfy the following coupled difference equations:
	\begin{align}\label{diffpn-1}
			 \textbf{p}(n)=\textbf{p}(n-1)-\mathit{i}\alpha_{n-1},
	\end{align}
\begin{equation}\label{diffbt-1}
\begin{aligned}
			\beta_{n+1}
=& \frac{\beta_{n-1}}{\eta\beta_{n}}\left[(2\gamma_{n}-3)\alpha_{n-1}+2(\gamma_{n}-2)\alpha_{n-2}\right.\\
&\left.\qquad\quad+\eta\left(\beta_{n-1}+\beta_{n-2}-\alpha_{n-1}^2-\alpha_{n-2}^2-\alpha_{n-2}\alpha_{n-1}-1\right)\right]\\
&+\frac{\alpha_{n-1}(\alpha_{n-1}^2+1) }{\eta\beta_{n}}+\frac{2i\textbf{p}(n-1)}{\eta}\left(\frac{\beta_{n-1}}{\beta_{n}}-1\right)+\alpha_{n}^2+\alpha_{n-1}^2+\alpha_{n}\alpha_{n-1}\\
&-\frac{1}{\eta}\left[2(\gamma_{n}+1)\alpha_{n}+(2\gamma_{n}+3)\alpha_{n-1}\right]-\beta_{n}+1,
\end{aligned}
\end{equation}
\begin{equation}\label{diffal-1}
\begin{aligned} \alpha_{n+1}=&\frac{1}{\eta\beta_{n+1}}\left[\left(\alpha_{n-1}-\alpha_{n}\right)\left(2\mathit{i}\textbf{p}(n)-\eta\right)+\left(\eta\alpha_{n}-2\gamma_{n}-3\right)\alpha_{n}^2-\left(\eta\alpha_{n-1}-2\gamma_{n}+3\right)\alpha_{n-1}^2\right.\\
&\left.\qquad\quad			 +\beta_{n}\big(\eta(\alpha_{n-1}-\alpha_{n})+2\big)+\beta_{n-1}\big(\eta (2\alpha_{n-1}+\alpha_{n-2})-2\gamma_{n}+3\big)-2\right]\\
&-2\alpha_{n}+\frac{2\gamma_{n}+3}{\eta},
\end{aligned}
\end{equation}	
for $n\geq3$, where $\gamma_{n}:=\gamma+n$.
	\end{proposition}
	\begin{proof}
Equation \eqref{diffpn-1} is actually  \eqref{al} with $n$ replaced by $n-1$.

Plugging \eqref{s2-4} into \eqref{(2.16)} and using \eqref{(3.1)} to eliminate $R_{n\pm1}$, we obtain an equation involving the recurrence coefficients and $ \textbf{p}(n\pm1)$. By replacing $\textbf{p}(n+1)$ with  $\textbf{p}(n)-\mathit{i}\alpha_{n}$ and $\textbf{p}(n-1)$ with $\mathit{i}\alpha_{n-1}+\textbf{p}(n)$ in this equation, we arrive at \eqref{diffbt-1}.

Replacing $n$ with $n-1$ in \eqref{(2.12)} gives
		\begin{align*}
			 r_{n}+r_{n-1}=1-\mathit{i}\alpha_{n-1}R_{n-1}.
		\end{align*}	
Substituting $\{r_{n},r_{n-1}\}$ and $R_{n-1}$ into the above identity using \eqref{(3.1)} and \eqref{(3.2)}, we get an equation involving the recurrence coefficients, $\textbf{p}(n)$ and $\textbf{p}(n-1)$. With $\textbf{p}(n-1)$ replaced by $\mathit{i}\alpha_{n-1}+\textbf{p}(n)$ in this equation, we come to \eqref{diffal-1}.
	\end{proof}

Rewriting \eqref{diffpn-1}-\eqref{diffal-1} as follows:
\begin{align*}
 \textbf{p}(n)=&\textbf{p}(n-1)-\mathit{i}\alpha_{n-1},\\
\beta_{n+1}=&f\left(\beta_n,\beta_{n-1},\beta_{n-2}, \alpha_n,\alpha_{n-1},\alpha_{n-2},\textbf{p}(n-1)\right),\\ \alpha_{n+1}=&g\left(\beta_{n+1},\beta_n, \beta_{n-1}, \alpha_n,\alpha_{n-1}, \alpha_{n-2}, \textbf{p}(n)\right),
 \end{align*}
for $n\geq3$, we find that they can be iterated in $n$, given the initial values $\textbf{p}(1)$ and $\{\alpha_i, \beta_i,i=1,2,3\}$. When $\lambda=0$ and $\zeta$ is a positive zero of the Bessel function $J_0$ , or $\lambda=1/2$ and $\zeta=m\pi$ with $m$ denoting a nonzero integer, we use \eqref{dp}-\eqref{al} and \eqref{bt} to calculate these initial values. By substituting them into \eqref{diffbt-1}-\eqref{diffal-1}  for $n=3,4,\dots$ and iterating, we can obtain $\alpha_{i},\beta_i$ and $\textbf{p}(i)$ for arbitrary $i$.

Eliminating $\textbf{p}(n)$ from \eqref{diffpn-1}-\eqref{diffal-1} leads us to the following theorem.
\begin{theorem}
The recurrence coefficients satisfy the following coupled difference equations:
\begin{equation}\label{diff-al}
\begin{aligned}
\alpha_{n+1}=&\frac{1}{\eta\beta_{n+1}}\left\{-3\alpha_{n}^2-2\alpha_{n}\alpha_{n-1}-\alpha_{n-1}^2+2(\beta_{n}-1)+\beta_{n-1}(\eta(2\alpha_{n-1}+\alpha_{n-2})-2\gamma_{n}+3)\right.\\
&\left.\qquad\quad\,+(\alpha_{n-1}-\alpha_{n})\big[(\alpha_{n-2}-\alpha_{n})\left(\eta(\alpha_{n}+\alpha_{n-1}+\alpha_{n-2})-2\gamma_n\right)+\eta(\beta_n-\beta_{n-1})\big]\right.\\
&\left.\qquad\quad\,+\frac{\alpha_{n-1}-\alpha_{n}}{\alpha_{n-2}-\alpha_{n-1}}\left[\alpha_{n-1}^2+5\alpha_{n-2}^2+2+\beta_{n}\left(\eta(\alpha_{n}+2\alpha_{n-1})-2\gamma_{n}-1\right)\right.\right.\\
&\left.\left.\quad\qquad\qquad\qquad\qquad\quad\,-2\beta_{n-1}+\beta_{n-2}\left(2\gamma_{n}-5-\eta(2\alpha_{n-2}+\alpha_{n-3})\right)\right]\right\}\\
&-2\alpha_{n}+\frac{2\gamma_{n}+3}{\eta},
\end{aligned}
\end{equation}
\begin{equation}\label{diff-bt}
	\begin{aligned}
\beta_{n+2}=&\alpha_{n+1}^2+\alpha_{n+1}\alpha_n+\alpha_{n}^2-\alpha_{n-1}^2-\beta_{n+1}+\beta_n+\beta_{n-1}-2(\gamma_{n}+2)\frac{\alpha_{n+1}}{\eta}-(2\gamma_{n}+5)\frac{\alpha_{n}}{\eta}\\
&+2(\gamma_n-1)\frac{\alpha_{n-1}}{\eta}+\frac{\alpha_{n}}{\eta\beta_{n+1}}\left[\alpha_n^2+1+\beta_{n}\left(2\gamma_{n}-1-\eta(\alpha_{n}+\alpha_{n-1})\right)\right]\\
&+\frac{\beta_{n+1}-\beta_{n}}{\eta\beta_{n+1}(\beta_{n-1}-\beta_{n})}\left\{\beta_{n}\left[\eta(\alpha_{n}^2+\alpha_n\alpha_{n-1}-\beta_{n+1})-2(\gamma_{n}+1)\alpha_{n}-3\alpha_{n-1}\right]\right.\\ &\left.\quad-\beta_{n-1}\left[\eta(\alpha_{n-2}^2+\alpha_{n-2}\alpha_{n-1}-\beta_{n-2})-2(\gamma_{n}-2)\alpha_{n-2}+3\alpha_{n-1}\right]+\alpha_{n-1}^3+\alpha_{n-1}\right\},
	\end{aligned}
\end{equation}
for $n\geq3$, where $\gamma_{n}:=n+\gamma$.
	\end{theorem}
	\begin{proof}
Solving for $\textbf{p}(n)$ from \eqref{diffal-1} yields
\begin{align*}
2\mathit{i}\textbf{p}(n)=&-2\gamma_n(\alpha_n+\alpha_{n-1})+\eta(\alpha_n^2+\alpha_n\alpha_{n-1}+\alpha_{n-1}^2-\beta_n+1)\\
&+\frac{1}{\alpha_{n-1}-\alpha_{n}}\left[3\left(\alpha_{n}^2+\alpha_{n-1}^2\right)-\beta_{n+1}\left(2\gamma_n+3-\eta(\alpha_{n+1}+2\alpha_{n})\right)\right.\\
&\left.\qquad\qquad\qquad\;-2\beta_n+\beta_{n-1}\left(2\gamma_n-3-\eta( 2\alpha_{n-1}+\alpha_{n-2})\right)+2\right].
\end{align*}
Replacing $n$ with $n-1$ in this expression, and subtracting the resulting identity from the above one, in view of $\mathit{i}(\textbf{p}(n)-\textbf{p}(n-1))=\alpha_{n-1}$, we come to \eqref{diff-al}.

Solving \eqref{diffbt-1} for $\textbf{p}(n-1)$  produces
\begin{equation*}
\begin{aligned}		 2\mathit{i}\textbf{p}(n-1)&=\frac{1}{\beta_{n}-\beta_{n-1}}\left\{\left[\eta(\alpha_{n}^2+\alpha_{n}\alpha_{n-1}-\beta_{n+1})-2(\gamma_n+1)\alpha_{n}-5\alpha_{n-1}\right]\beta_{n}\right.\\ &\left.+\left[-\eta(\alpha_{n-2}^2+\alpha_{n-1}\alpha_{n-2}-\beta_{n-2})-\alpha_{n-1}+2(\gamma_n-2)\alpha_{n-2}\right]\beta_{n-1}+\alpha_{n-1}^3+\alpha_{n-1}\right\}\\ &+\eta(\alpha_{n-1}^2-\beta_{n}-\beta_{n-1}+1)+2(1-\gamma_n)\alpha_{n-1}.\label{(1.14)}
	\end{aligned}
\end{equation*}
Substituting $n+1$ for $n$ in the above expression and subtracting  the resulting identity from the above one,  in light of $\mathit{i}(\textbf{p}(n-1)-\textbf{p}(n))=-\alpha_{n-1}$, we arrive at \eqref{diff-bt}.
\end{proof}

Rewriting \eqref{diff-al} and \eqref{diff-bt} as follows:
\begin{align*}
\alpha_{n+1}=&G\left(\beta_{n+1},\beta_{n}, \beta_{n-1}, \beta_{n-2}, \alpha_{n},\alpha_{n-1}, \alpha_{n-2}, \alpha_{n-3}\right),\\
\beta_{n+2}=&F\left(\beta_{n+1},\beta_{n},\beta_{n-1},\beta_{n-2}, \alpha_{n+1},\alpha_{n},\alpha_{n-1},\alpha_{n-2}\right), \end{align*}
for $n\geq3$, we find that they can be iterated in $n$, given the initial values $\{\alpha_0,\alpha_1,\alpha_2,\alpha_3\}$ and $\{ \beta_1,\beta_2,\beta_3,\beta_4\}$. When $\gamma=0$ and $\eta$ is a positive zero of the Bessel function $J_0$ , or $\gamma=1/2$ and $\eta=m\pi$ with $m$ denoting a nonzero integer, we use \eqref{dp}-\eqref{al} and \eqref{bt} to calculate these initial values. By substituting them into \eqref{diff-al} and \eqref{diff-bt} for $n=3,4,\dots$ and iterating, we obtain $\alpha_{i}$ and $\beta_i$ for arbitrary $i$. In particular, for $n=3,4,5,6,7,8$, we get the expressions for $\alpha_i (i=4,\dots,9)$ and $\beta_i (i=5,\dots,10)$ in $\zeta$, which together with the initial conditions are presented in Tables \ref{tab1:rec-coeffs} and \ref{tab2:rec-coeffs}. Note that $\beta_{10}$ is not included in the tables.

\begin{corollary}\label{cor0}
When $\gamma=0$ and $\eta$ is a positive zero of the Bessel function $J_0$, the expressions of $\{\alpha_{n}, \beta_n,0\leq n\leq9\}$ are displayed in Table \ref{tab1:rec-coeffs}.
		\begin{table}[H]
			\centering
			 \renewcommand{\arraystretch}{1.5}
			 \setlength{\tabcolsep}{3pt}
\caption{Recurrence coefficients $\alpha_{n}$ and $\beta_{n}$ for $0\leq n\leq9$}
			\label{tab1:rec-coeffs}
\begin{tabular}{|c|c|c|}
\hline
$n$ & $\alpha_n$ & $\beta_n$   \\
\hline
0 & $\frac{1}{\eta}$ & $\beta_0$ \\
\hline
1 & $\frac{3}{\eta}+\frac{-3\eta}{\eta^2-1}$ & $1-\frac{1}{\eta^2}$  \\
\hline
2 & $\frac{5}{\eta}+3\eta\left(\frac{1}{\eta^2-1}-\frac{\eta^2-14}{2W_2}\right)$ & $\frac{-2W_2}{\eta^2(\eta^2-1)^2}$ \\
\hline
3 & $\frac{7}{\eta}+\eta\left(\frac{3(\eta^2-14)}{2W_2}-10\frac{ W_3}{W_4}\right)$ & $\frac{(\eta^2-1)W_4}{\eta^2W_2^2}$ \\
\hline
4 & $\frac{9}{\eta}+10\eta\left(\frac{W_3}{W_4}-\frac{W_5}{2W_6}\right)$ & $-12\frac{W_2W_6}{\eta^2W_4^2}$ \\
\hline
5 & $\frac{11}{\eta}+\eta\left(5\frac{W_5}{W_6}-21\frac{W_{8}}{W_{9}}\right)$ & $\frac{W_4W_9}{\eta^2W_6^2}$ \\
\hline
6 & $\frac{13}{\eta}+21\eta\left(\frac{W_{8}}{W_{9}}-\frac{W_{11}}{2W_{12}}\right)$ & $-30\frac{W_6W_{12}}{\eta^2W_9^2}$ \\
\hline
7 & $\frac{15}{\eta}+6\eta\left(7\frac{W_{11}}{4W_{12}}-6\frac{W_{15}}{W_{16}}\right)$ & $\frac{W_9W_{16}}{\eta^2W_{12}^2}$ \\
\hline
8&$\frac{17}{\eta}+36\eta\left(\frac{W_{15}}{W_{16}}-\frac{W_{19}}{2W_{20}}\right)$&$-56\frac{W_{12}W_{20}}{\eta^2W_{16}^2}$\\
\hline
9&$\frac{19}{\eta}+\eta\left(18\frac{W_{19}}{W_{20}}-55\frac{W_{24}}{W_{25}}\right)$&$\frac{W_{16}W_{25}}{\eta^2W_{20}^2}$\\
\hline
\end{tabular}
\par {\small Note: $W_i=W_i(\eta^2)$, with $W_i(\cdot)$ denoting a monic polynomial of degree $i$. Hence,\\ $W_i(\eta^2)$ is a monic polynomial of degree $2i$ in $\eta$, containing only even powers of $\eta$.\\ The expressions for $W_i(\eta^2), i=2,3,4,5,6,8,9,11,12,15,16$, are given in \\Appendix \ref{symbols}, while the others are omitted for brevity.}
\end{table}
		
\end{corollary}

\begin{proof}
When $\gamma=0$ in \eqref{w2}, the weight function now reads
\begin{align}
w_2(x)=x(1-x^2)^{-\frac{1}{2}}\exp(\mathit{i}\eta x).
\end{align}
Given the initial values
\[P_{-1}(x)=0,\qquad P_0(x)=1,\]
we use the integral expressions \eqref{al-1} and \eqref{bt} to compute $\alpha_n$ and $\beta_n$ respectively, i.e.
\begin{gather}
\alpha_n=\frac{-i}{h_n}\int_{-1}^1xP_n^2(x)w(x)dx,\label{pf0-eq2}\\
\beta_n=\frac{h_n}{h_{n-1}},
\end{gather}
where $h_n$ is given by \eqref{dp} with $m=n$ and $P_n$ satisfies the recurrence relation \eqref{3r}, namely,
\begin{gather}
h_n=\int_{-1}^1 P_n^2(x)w(x)dx,\label{pf0-eq1}\\
P_{n+1}(x)=(x-i\alpha_{n})P_{n}(x)-\beta_{n}P_{n-1}(x),\label{pf0-eq5}
\end{gather}
for $n\geq0$. Using \eqref{pf0-eq2}-\eqref{pf0-eq5}, with the Mathematica code provided in Appendix \ref{A4}, we get the values of $\alpha_i$ for $i=0,1,2,3$ and $\beta_i$ for $i=1,2,3,4$.
With these initial values, by iterating the difference equations \eqref{diffpn-1}-\eqref{diffal-1} for $n=3,4,5,6,7,8$, we generate the values of $\alpha_n$ and $\beta_n$ listed in Table \ref{tab1:rec-coeffs}.
\end{proof}


\begin{remark}
The initial values $\{\alpha_0,\alpha_1,\alpha_2,\alpha_3\}$ and $\{\beta_1,\beta_2,\beta_3,\beta_4\}$ can be derived by hand. The derivation process is provided in Appendix \ref{J0D-int}.
\end{remark}

From the expressions of $\alpha_n$ and $\beta_n$ presented in Table \ref{tab1:rec-coeffs}, we propose the following conjecture.
\begin{conjecture} The recurrence coefficients for the monic orthogonal polynomials associated with the weight function $x(1-x^2)^{-1/2}\exp(i\eta x), x\in[-1,1]$, with $\zeta$ being a positive zero of $J_{0}$, have the following symbolic forms:
\begin{align*}
\alpha_{2k}=&\frac{4k+1}{\eta}+k(2k+1)\eta\left(\frac{W_{k^2-1}(\eta^2)}{W_{k^2}(\eta^2)}-\frac{W_{k(k+1)-1}(\eta^2)}{2W_{k(k+1)}(\eta^2)}\right),\\ \alpha_{2k+1}=&\frac{4k+3}{\eta}+\eta\left(k(2k+1)\frac{W_{k(k+1)-1}(\eta^2)}{2W_{k(k+1)}(\eta^2)}-(k+1)(2k+3)\frac{W_{k(k+2)}(\eta^2)}{W_{(k+1)^2}(\eta^2)}\right),
\end{align*}
for $k\geq2$, and
\begin{align*}
\beta_{2k} =& -2k(2k-1)\frac{W_{k(k-1)}(\eta^2)\cdot W_{k(k+1)}(\eta^2)}{\eta^2 \left(W_{k^2}(\eta^2)\right)^2},& \beta_{2k+1} =& \frac{W_{k^2}(\eta^2)\cdot W_{(k+1)^2}(\eta^2)}{\eta^2 \left(W_{k(k+1)}(\eta^2)\right)^2},
\end{align*}
for $k\geq1$. Here $W_j(\cdot)$ is a monic polynomial of degree $j$ for $j\geq0$, with $W_0(\eta^2):=1$ and $W_1(\eta^2):=\eta^2-1$. Note that $W_i(\eta^2)$ for $i=2,3,4,5,6,8,9,11,12,15,16$ are listed in Appendix \ref{symbols}.
\end{conjecture}

When $\gamma=\frac{1}{2}$ and $\eta=m\pi$ where $m$ is an arbitrary nonzero integer, the weight function now reads
\begin{align}
w_2(x)=x\exp(\mathit{i}\eta x).
\end{align}	
Via an argument similar to that in the proof of Corollary \ref{cor0}, we generate the values of $\alpha_n$ and $\beta_n$ listed in Table \ref{tab2:rec-coeffs}. It should be pointed out that, the initial values can be derived using Mathematica code similar to that in Appendix \ref{A4} (see the comments therein for details).
\begin{corollary}
When $\gamma=\frac{1}{2}$ and $\eta=m\pi$ where $m$ is an arbitrary nonzero integer, the values of $\{\alpha_{n},\beta_n, 0\leq n\leq 9\}$ are given in Table \ref{tab2:rec-coeffs}.	
		\begin{table}[H]
			\centering
			 \renewcommand{\arraystretch}{1.5}
			 \setlength{\tabcolsep}{3pt}
			\caption{Recurrence coefficients $\alpha_{n}$ and $\beta_{n}$ for $0\leq n \le 9$}
			\label{tab2:rec-coeffs}
			 \begin{tabular}{|c|c|c|}
\hline
$n$ & $\alpha_n$ & $\beta_n$   \\
\hline
0 & $\frac{2}{\eta}$ & $\beta_0$  \\
\hline
1 & $\frac{4}{\eta}-\frac{4\eta}{\eta^2-2}$ & $1-\frac{2}{\eta^2}$ \\
\hline
2&$\frac{6}{\eta}+2\eta\left(\frac{2}{\eta^2-2}-\frac{\eta^2-24}{G_2}\right)$ & $\frac{-4G_2}{\eta^2(\eta^2-2)^2}$ \\
\hline
3&$\frac{8}{\eta}+2\eta\left(\frac{\eta^2-24}{G_2}-6\frac{G_3}{G_4}\right)$&$\frac{(\eta^2-2)G_4}{\eta^2G_2^2}$ \\
\hline	
4&$\frac{10}{\eta}+6\eta\left(2\frac{G_3}{G_4}-\frac{G_5}{G_6}\right)$ & $-16\frac{G_2G_6}{\eta^2G_4^2}$ \\
\hline
5&$\frac{12}{\eta}+6\eta\left(\frac{G_5}{G_6}-\frac{4G_8}{G_9}\right)$ & $\frac{G_4G_9}{\eta^2G_6^2}$ \\
\hline
6&$\frac{14}{\eta}+12\eta\left(\frac{2G_8}{G_9}-\frac{G_{11}}{G_{12}}\right)$ & $-36\frac{G_6G_{12}}{\eta^2G_9^2}$ \\
\hline
7&$\frac{16}{\eta}+2\eta\left(\frac{6G_{11}}{G_{12}}-20\frac{G_{15}}{G_{16}}\right)$ & $\frac{G_9G_{16}}{\eta^2G_{12}^2}$ \\\hline
8&$\frac{18}{\eta}+20\eta\left(\frac{2G_{15}}{G_{16}}-\frac{G_{19}}{G_{20}}\right)$&$-64\frac{G_{12}G_{20}}{\eta^2G_{16}^2}$\\
\hline
9&$\frac{20}{\eta}+20\eta\left(\frac{G_{19}}{G_{20}}-3\frac{G_{24}}{G_{25}}\right)$&$\frac{G_{16}G_{25}}{\eta^2G_{20}^2}$\\
\hline
\end{tabular}
\par {\small Note: $G_i=G_i(\eta^2)$, with $G_i(\cdot)$ denoting a monic polynomial of degree $i$. Hence, \\ $G_i(\eta^2)$ is a monic polynomial of degree $2i$ in $\eta$, containing only even powers of $\eta$.\\ The expressions for $G_i(\eta^2), i=2,3,4,5,6,8,9,11,12,15,16,$ are given in\\ Appendix \ref{symbols}, while the others are omitted for brevity.}
\end{table}
\end{corollary}

According to the above expressions of $\alpha_n$ and $\beta_n$, we propose the following conjecture.
\begin{conjecture} The recurrence coefficients for the monic orthogonal polynomials associated with the weight function $x\exp(im\pi x), x\in[-1,1]$, with $m$ being an arbitrary nonzero integer, have the following symbolic forms:
\begin{align*}
\alpha_{2k}=&\frac{2(2k+1)}{\eta}+k(k+1)\eta\left(\frac{2G_{k^2-1}(\eta^2)}{G_{k^2}(\eta^2)}-\frac{G_{k(k+1)-1}(\eta^2)}{G_{k(k+1)}(\eta^2)}\right),\\ \alpha_{2k+1}=&\frac{2(2k+2)}{\eta}+\eta\left(k(k+1)\frac{G_{k(k+1)-1}(\eta^2)}{G_{k(k+1)}(\eta^2)}-2(k+1)(k+2)\frac{G_{k(k+2)}(\eta^2)}{G_{(k+1)^2}(\eta^2)}\right),
\end{align*}
for $k\geq2$, and
\begin{align*}
\beta_{2k} =& -(2k)^2\frac{G_{k(k-1)}(\eta^2)\cdot G_{k(k+1)}(\eta^2)}{\eta^2 \left(G_{k^2}(\eta^2)\right)^2},& \beta_{2k+1} =& \frac{G_{k^2}(\eta^2)\cdot G_{(k+1)^2}(\eta^2)}{\eta^2 \left(G_{k(k+1)}(\eta^2)\right)^2},
\end{align*}
for $k\geq1$. Here $G_j(\cdot)$ is a monic polynomial of degree $j$ for $j\geq0$, with $G_0(\eta^2):=1$ and $G_1(\eta^2):=\eta^2-2$. Note that $G_i(\eta^2)$ for $i=2,3,4,5,6,8,9,11,12,15,16$ are listed in Appendix \ref{symbols}.
\end{conjecture}

\begin{remark}
Our conjecture is consistent with the one given in \cite[Conjecture 3.1]{MC05-2}.
\end{remark}



\begin{remark}
The values of $\{\alpha_0,\alpha_1,\alpha_2,\alpha_3,\alpha_4;\beta_1,\beta_2,\beta_3,\beta_4\}$ given in Table \ref{tab2:rec-coeffs} can be calculated by hand. In fact, for $\eta=m\pi$, using Euler's formula $e^{i\eta x}=\cos (\eta x)+i \sin (\eta x)$ and through integration by parts, we obtain the following integral identities:
\begin{align*}
\int_{-1}^1 x^k e^{i\eta x}dx=&\frac{-2i(-1)^m}{\eta^k}\tau_k,\qquad 1\leq k\leq 10,
\end{align*}
where $\tau_1=1, \tau_2=2i,\tau_3=\eta^2-6,\tau_4=4i\tau_3, \tau_5=\eta^4-20 \eta^2+120,\tau_6=6i\tau_5,\tau_7=\eta ^6-42
   \eta ^4+840 \eta
   ^2-5040,\tau_8=8i\tau_7, \tau_9=\eta ^8-72 \eta ^6+3024
   \eta ^4-60480 \eta
   ^2+362880, \tau_{10}=10i\tau_9$.
Using \eqref{pf0-eq2}-\eqref{pf0-eq5}, with the aid of the above integral formulas, we obtain the desired values.
\end{remark}





\begin{remark}

When $\gamma=\frac{1}{2}$ and $\eta=100\pi$, with the aid of the initial values $\{\alpha_0,\alpha_1,\alpha_2,\alpha_3\}$ and $\{\beta_1,\beta_2,\beta_3,\beta_4\}$ given in Table \ref{tab2:rec-coeffs}, we iterate the coupled difference equations \eqref{diff-al}-\eqref{diff-bt} for $3\leq n\leq 30$. The Mathematica code is provided in Appendix \ref{A6}, and the obtained values of $\alpha_n$ and $\beta_n$ for $0\leq n \leq 29$ are the same as those given in \cite[Table 2]{MC05-2} which were computed using a combination of the Chebyshev method and the Stieltjes-Gautschi procedure. In addition, we get
\begin{align*}
\alpha_{30}=&1.269928185955803899,&
\beta_{30}=&-0.00428212384856920060.
\end{align*}


For $\gamma=1/2$ and $\eta=\pi$, using the Mathematica code similar to that given in Appendix \ref{A6} (see the comments therein), we compute $\alpha_n$ and $\beta_n$ for $0\leq n \leq 29$. The obtained values are identical to those in \cite[Table 2]{MC05-2}, except that the last digit of $\alpha_{18}$ is 3 (not 2), the last digit of $\beta_8$ is 8 (not 7), and the last two digits of $\beta_{22}$ are 70 (not 69). In addition, we get
\begin{align*}
\alpha_{30}=&-0.08657434839001662,&
\beta_{30}=&0.2103637730884711.
\end{align*}
\end{remark}

\begin{remark}
Since $\eta$ is real in both cases $\gamma=0$ and $\gamma=1/2$,  we know that the initial values $\{\alpha_i,\beta_i,i=1,2,3,4\}$ given in Table \ref{tab1:rec-coeffs} and Table \ref{tab2:rec-coeffs} are all rational numbers. Noting that the coefficients that appear in the difference equations \eqref{diffpn-1}-\eqref{diffal-1} are all rational, we conclude that the values of $\alpha_n (n\geq0)$ and $\beta_n (n\geq1)$ are rational.
\end{remark}


\section*{Acknowledgements}
This work was supported by National Natural Science Foundation of China under grant numbers 12101343 and 12371257, and by Shandong Provincial Natural Science Foundation with project number ZR2021QA061.




\bibliographystyle{}
\begin{thebibliography}{10}

\bibitem{BCH}
E. Basor, Y. Chen and N. Haq, Asymptotics of determinants of Hankel matrices via non-linear difference equations, J. Approx. Theory {\bf 198} (2015), 63-110.

\bibitem{CI}
Y. Chen and M. Ismail, Jacobi polynomials from compatibility conditions, Proc. Am. Math. Soc. {\bf 133} (2005), 465-472.

\bibitem{CI10}
Y. Chen and A. Its, Painlev\'{e} III and a singular linear statistics in Hermitian random matrix ensembles, I, J. Approx. Theory {\bf 162} (2010), 270-297.


\bibitem{CM12}
Y. Chen and M. McKay, Coulomb fluid, Painlev\'{e} transcendents and the information theory of MIMO systems, IEEE Trans. Inf. Theory {\bf 58} (2012), 4594-4634.

\bibitem{CYC}
R. Chen, D. Yu and J. Chen, Computation of oscillatory integrals with an exponential kernel
and Jacobi-type singularities, J. Math. Anal. Appl. {\bf 494} (2021), 124448 (11pp).

\bibitem{CJ21}
P. A. Clarkson and K. Jordaan, Generalised Airy polynomials, J. Phys. A: Math. Theor. {\bf 54} (2021), 185202 (28pp).

\bibitem{GR}
I. S. Gradshteyn and I. M. Ryzhik, Table of Integrals, Series, and Products, Seventh edition. Edited by Daniel Zwillinger and Victor Moll. Academic Press, Elsevier, 2015.

\bibitem{Ismail}
M. Ismail, Classical and Quantum Orthogonal Polynomials in One Variable, Encyclopedia of Mathematics and its Applications 98, Cambridge University Press, Cambridge, 2005.

\bibitem{KX}
H. Kang and Q. Xu, Quadrature formulae of many highly oscillatory Frourier-type integrals with algebraic or logarithmic singularities and their error analysis, Appl. Math. Comput. {\bf 442} (2023), 127758 (15pp).

\bibitem{KHM}
D. K. Kurto\v{g}lu, A. I. Has\c{c}elik and G. V. Milovanovi\'c, A method for efficient computation of integrals with oscillatory and singular integrand, Numer. Algorithms {\bf 85} (2020), 1155-1173.

\bibitem{Lebedev}
N. N. Lebedev, Special Functions and Their Applications, Revised English Edition Translated and Edited by Richard A. Silverman. Dover Publications, New York, 1972.


\bibitem{LyuLyu25}
S. Lyu and Y. Lyu, Ladder Operators for Laguerre-type and Jacobi-type Orthogonal Polynomials, J. Phys. A: Math. Theor. {\bf 59} (2026), 115201 (30pp).

\bibitem{MC05-1}
G. V. Milovanovi\'c and A. S. Cvetkovi\'c,  Orthogonal polynomials related to the oscillatory Chebyshev weight function, {Bull. Cl. Sci. Math. Nat. Sci. Math. 30 (2005)}, {47-60}.


\bibitem{MC05-2}
G V. Milovanovi\'c and A. S. Cvetkovi\'c, Orthogonal polynomials and Gaussian quadrature rules related to oscillatory weight functions, J. Comput. Appl. Math. 179 (2005), 263-287.


\bibitem{MCM}
G. V. Milovanovi\'{c}, A. S. Cvetkovi\'{c} and Z. M. Marjanovic, Orthogonal polynomials for the oscillatory-Gegenbauer weight, Publ. Inst. Math. (Belgr.) (N.S.) {\bf 84} (2008), 49-60.


\bibitem{MCS}
G. V. Milovanovi\'{c}, A. S. Cvetkovi\'{c} and M. P. Stani\'{c}, Orthogonal polynomials for modified Gegenbauer weight and corresponding quadratures, Appl. Math. Lett. {\bf 22} (2009), 1189-1194.


\bibitem{MSM}
G. V. Milovanovi\'{c}, M. P. Stani\'{c} and T. V. Tomovi\'{c} Mladenovi\'{c}, Gaussian type quadrature rules related to the oscillatory modification of the generalized Laguerre weight functions, J. Comput. Appl. Math. {\bf 437} (2024), 115476 (8pp).

\bibitem{MC17}
C. Min and Y. Chen, Gap probability distribution of the Jacobi unitary ensemble: an elementary treatment, from finite $n$ to double scaling, Stud. Appl. Math. {\bf 140} (2017), 202-220.

\bibitem{MC20}
C. Min and Y. Chen, Painlev\'{e} V and the Hankel determinant for a singularly perturbed Jacobi weight, Nucl. Phys. B {\bf 961} (2020), 115221 (25pp).

\bibitem{MC21}
C. Min and Y. Chen, Differential, difference and asymptotic relations for Pollaczek-Jacobi type orthogonal polynomials and their Hankel determinants, Stud. Appl. Math. {\bf 147} (2021), 390-416.


\bibitem{MF25}
C. Min and P. Fang, Generalized Airy polynomials, Hankel determinants and asymptotics, Physica D {\bf 473} (2025), 134560 (9pp).


\bibitem{MW26-1}
C. Min and X. Wu, Asymptotics of the Hankel determinant and orthogonal polynomials arising from the information theory of MIMO systems, arXiv: 2510.06739.


\bibitem{MW26-2}
C. Min and X. Wu, Orthogonal polynomials for the singularly perturbed Laguerre weight, Hankel determinants and asymptotics, arXiv: 2511.09362.


\bibitem{ML25}
X. Mu and S. Lyu, Hankel determinants for a deformed Laguerre weight with multiple variables and generalized Painlev\'{e} V equation, J. Math. Phys. {\bf 66} (2025), 073506 (21pp).


\bibitem{SC}
M. P. Stani\'{c} and A. S. Cvetkovi\'{c}, Orthogonal polynomials with respect to modified
Jacobi weight and corresponding quadrature rules
of Gaussian type, Numer. Math. Theor. Meth. Appl. {\bf 4} (2011), 478-488.


\end{thebibliography}


{\small \begin{appendices}

\section{Derivation of the initial values $\{\alpha_0,\alpha_1\}$ and $\{ \beta_1,\beta_2\}$ given by \eqref{al0bt1}.}\label{J0D-int-1}
We need the following lemma.
\begin{lemma}\label{w1L}
Supposing $\zeta$ is a positive zero of the Bessel function $J_{\lambda-1}$, we have the following six integral identities:
\begin{align}
\int_{-1}^{1}(1-x^2)^{\lambda-\frac{1}{2}}\exp(\mathit{i}\zeta x)\, dx=&A\zeta^{-\lambda}J_{\lambda}(\zeta),\label{Bs-1-1}\\
\int_{-1}^{1}x(1-x^2)^{\lambda-\frac{1}{2}}\exp(\mathit{i}\zeta x)\, dx=&2i A\lambda\zeta^{-\lambda-1}J_{\lambda}(\zeta),\label{Bs-1-2}\\
\int_{-1}^{1}x^{2}(1-x^2)^{\lambda-\frac{1}{2}}\exp(\mathit{i}\zeta x)\, dx=&A\zeta^{-\lambda}J_{\lambda}(\zeta)\left[1-\frac{2\lambda(2\lambda+1)}{\zeta^2}\right],\label{Bs-1-3}\\
\int_{-1}^{1}x^{3}(1-x^2)^{\lambda-\frac{1}{2}}\exp(\mathit{i}\zeta x)\, dx=&iA\zeta^{-\lambda-1}J_{\lambda}(\zeta)\left[4\lambda+1-\frac{2\lambda(2\lambda+1)(2\lambda+2)}{\zeta^2}\right],\label{Bs-1-4}\\
\int_{-1}^{1}x^{4}(1-x^2)^{\lambda-\frac{1}{2}}\exp(\mathit{i}\zeta x)\, dx=&A\zeta^{-\lambda}J_{\lambda}(\zeta)\left[1-\frac{3(2\lambda+1)^2}{\zeta^2}+\frac{2\lambda(2\lambda+1)(2\lambda+2)(2\lambda+3)}{\zeta^4}\right],\label{Bs-1-5}
\end{align}
where $A:=\sqrt{\pi}2^{\lambda}\Gamma\left(\lambda+\frac{1}{2}\right)$.
\end{lemma}

\begin{proof}
We first recall the following properties of the Bessel function \cite[(5.3.6)-(5.3.7)]{Lebedev}:
\begin{align}
\frac{2\nu}{z}J_{\nu}(z)=&J_{\nu-1}(z)+J_{\nu+1}(z),\label{Bs1}\\
2J_{\nu}'(z)=&J_{\nu-1}(z)-J_{\nu+1}(z).\label{Bs2}
%\frac{d}{dz}\left(z^{-\nu}J_{\nu}(z)\right)=&-z^{-\nu}J_{\nu+1}(z).\label{Bs3}
\end{align}
Solving $J_{\nu+1}(z)$ from \eqref{Bs1} and \eqref{Bs2} yields
\begin{align}\label{Bs3}
J_{\nu}'(z)=\frac{\nu}{z}J_{\nu}(z)-J_{\nu+1}(z).
\end{align}
From \eqref{Bs2} follows
\begin{align}
J_{\nu+1}'(z)=&\frac{1}{2}\left(J_{\nu}(z)-J_{\nu+2}(z)\right)\nonumber\\
=&J_{\nu}(z)-\frac{\nu+1}{z}J_{\nu+1}(z),\label{Bs4}
\end{align}
where the second identity is due to the following relation which comes from \eqref{Bs1}:
\[J_{\nu+2}(z)=\frac{2(\nu+1)}{z}J_{\nu+1}(z)-J_{\nu}(z).\]
Noting that $J_{\lambda-1}(\zeta)=0$, we get from \eqref{Bs1} that
\begin{align}\label{Bs5}
\frac{2\lambda}{\zeta}J_{\lambda}(\zeta)=J_{\lambda+1}(\zeta).
\end{align}

To continue, from \cite[eq. 10, p. 912]{GR}, we know that
\begin{align}\label{h0-1}
\int_{-1}^1(1-x^2)^{\nu-\frac{1}{2}}\exp(iz x)dx=Az^{-\nu}J_{\nu}(z),
\end{align}
which gives \eqref{Bs-1-1}.
Differentiating it with respect to $z$, with the help of \eqref{Bs3}, we find
\begin{align}
i\int_{-1}^{1}x(1-x^2)^{\nu-\frac{1}{2}}\exp(\mathit{i}z x)\, dx
=&-Az^{-\nu}J_{\nu+1}(z),\label{x1-2}
\end{align}
which, in light of \eqref{Bs5}, leads us to \eqref{Bs-1-2}.
By a subsequent differentiation of \eqref{x1-2} with respect to $z$, with the aid of \eqref{Bs4}, we obtain
\begin{align}
-\int_{-1}^{1}x^2(1-x^2)^{\nu-\frac{1}{2}}\exp(\mathit{i}z x)\, dx
=&- Az^{-\nu}\left(J_{\nu}(z)-\frac{2\nu+1}{z} J_{\nu+1}(z)\right),\label{Bs6}
\end{align}
which, in view of \eqref{Bs5}, gives us \eqref{Bs-1-3}.

Taking the derivative of \eqref{Bs6} with respect to $z$ and applying \eqref{Bs3}-\eqref{Bs4} results in
\begin{align}
i\int_{-1}^{1}x^3(1-x^2)^{\nu-\frac{1}{2}}\exp(\mathit{i}z x)\, dx
=& Az^{-\nu}\left[-\frac{2\nu+1}{z}J_{\nu}(z)+\left(\frac{(2\nu+1)(2\nu+2)}{z^2}-1\right)J_{\nu+1}(z)\right].\label{Bs7}
\end{align}
According to \eqref{Bs5}, we come to \eqref{Bs-1-4}. Finally, differentiating \eqref{Bs7} once more, we deduce
\begin{align*}
-\int_{-1}^{1}x^4(1-x^2)^{\nu-\frac{1}{2}}\exp(\mathit{i}z x)\, dx
=& Az^{-\nu}\left[\left(\frac{(2\nu+1)(2\nu+3)}{z^2}-1\right)J_{\nu}(z)\right.\\
&\left.\qquad\qquad+\left(-\frac{(2\nu+2)(2\nu+3)}{z^3}+\frac{2}{z}\right)(2\nu+1)J_{\nu+1}(z)\right],
\end{align*}
which, by virtue of \eqref{Bs5} again, leads to \eqref{Bs-1-5}.
\end{proof}

Recall \eqref{al-1} and \eqref{bt}:
\begin{gather}
\alpha_n=\frac{-i}{h_n}\int_{-1}^1xP_n^2(x)(1-x^2)^{\lambda-\frac{1}{2}}\exp(\mathit{i}\zeta x)dx, \qquad\qquad n\geq0,\label{al-int}\\
\beta_n=\frac{h_n}{h_{n-1}}, \qquad\qquad n\geq1,\label{bt-h}
\end{gather}
where $h_n$ is given by \eqref{dp} with $m=n$, namely
\begin{gather}\label{h-int}
h_n=\int_{-1}^1P_n^2(x)(1-x^2)^{\lambda-\frac{1}{2}}\exp(\mathit{i}\zeta x)dx,\qquad\qquad n\geq0,
\end{gather}
and $P_n$ can be computed using the three-term recurrence relation, i.e.
\begin{gather}
P_{n+1}(x)=(x-i\alpha_{n})P_{n}(x)-\beta_{n}P_{n-1}(x),\qquad\qquad n\geq0,\label{P-3r}
\end{gather}
with $P_{-1}(x):=0$ and $P_0(x):=1$.

With $n=0$ in \eqref{h-int}, \eqref{al-int} and \eqref{P-3r}, applying \eqref{Bs-1-1} and \eqref{Bs-1-2}, we have
\begin{align*}
h_0=&\int_{-1}^1(1-x^2)^{\lambda-\frac{1}{2}}\exp(\mathit{i}\zeta x)dx=A\zeta^{-\lambda}J_{\lambda}(\zeta),\\
\alpha_0=&\frac{-i}{h_0}\int_{-1}^1x(1-x^2)^{\lambda-\frac{1}{2}}\exp(\mathit{i}\zeta x)dx=\frac{2\lambda}{\zeta},\\
P_1(x)=&(x-i\alpha_{0})P_{0}(x)-\beta_{0}P_{-1}(x)=x-\frac{2 i\lambda}{\zeta}.
\end{align*}
Setting $n=1$ in \eqref{h-int}, \eqref{al-int}-\eqref{bt-h} and \eqref{P-3r}, by applying \eqref{Bs-1-1}-\eqref{Bs-1-4}, we get
\begin{align*}
h_1=&\int_{-1}^1P_1^2(x)(1-x^2)^{\lambda-\frac{1}{2}}\exp(\mathit{i}\zeta x)dx=A\zeta^{-\lambda}J_{\lambda}(\zeta)\left(1-\frac{2\lambda}{\zeta^2}\right),\\
\alpha_1=&\frac{-i}{h_1}\int_{-1}^1xP_1^2(x)(1-x^2)^{\lambda-\frac{1}{2}}\exp(\mathit{i}\zeta x)dx=\frac{\zeta^2-4\lambda(\lambda+1)}{\zeta(\zeta^2-2\lambda)},\\
\beta_1=&\frac{h_1}{h_0}=1-\frac{2\lambda}{\zeta^2},\\
P_{2}(x)=&(x-i\alpha_{1})P_{1}(x)-\beta_{1}P_{0}(x)=x^2-\frac{i (2\lambda +1) \left(\zeta^2-4 \lambda \right)}{\zeta(\zeta^2-2 \lambda) } x +\frac{4 \lambda ^2 (2 \lambda+1)}{\zeta^2(\zeta^2-2\lambda)}-1.
\end{align*}
Substituting $P_2(x)$ into \eqref{h-int} with $n=2$, with the aid of \eqref{Bs-1-1}-\eqref{Bs-1-5}, we obtain
\[h_2=\int_{-1}^1P_2^2(x)(1-x^2)^{\lambda-\frac{1}{2}}\exp(\mathit{i}\zeta x)dx=-\frac{2(2 \lambda +1) A \zeta ^{-\lambda}J_{\lambda}(\zeta)
   \left(\zeta ^4-\zeta ^2
   \lambda  (2 \lambda +5)+4
   \lambda ^2\right)}{\zeta^4(\zeta
   ^2-2 \lambda) },\]
so that
\[\beta_2=\frac{h_2}{h_1}=-\frac{2 (2 \lambda +1)\left(\zeta ^4-\zeta ^2\lambda  (2 \lambda +5)+4\lambda^2\right)}{\zeta^2\left(\zeta ^2-2\lambda \right)^2}.\]

\section{Mathematica code to compute $\{\alpha_0,\alpha_1\}$ and $\{ \beta_1,\beta_2\}$ given by \eqref{al0bt1}.}\label{IV1-M}
\begin{lstlisting}[language=Mathematica]
(*initialization*)
P[0, x_] = 1;
P[-1, x_] = 0;
w[x_] = (1 - x^2)^(\[Lambda] - 1/2)*Exp[I*\[Zeta] x];
$Assumptions = Element[\[Lambda], Rationals] && \[Lambda] > -1/2;

(*definitions of h[0] and \[Alpha][0]*)
h[0] = FullSimplify[Integrate[w[x], {x, -1, 1}]];
\[Alpha][0] =
  FullSimplify[
   FullSimplify[-I/h[0]*Integrate[x*w[x], {x, -1, 1}]] /.
    BesselJ[\[Lambda] +1, \[Zeta]] -> (2 \[Lambda]/\[Zeta]) BesselJ[\[Lambda], \[Zeta]]];

(*calculations of \[Alpha][n] and \[Beta][n] for n=1,2*)
results =
  Table[If[n > 0,
    P[n,x_] = (x - I*\[Alpha][n - 1])*P[n - 1, x] - \[Beta][n - 1]* P[n - 2, x];
    h[n] =
     FullSimplify[
      FullSimplify[Integrate[P[n, x]^2*w[x], {x, -1, 1}]] /. BesselJ[\[Lambda]+1, \[Zeta]] -> (2 \[Lambda]/\[Zeta]) BesselJ[\[Lambda], \[Zeta]]];
    \[Beta][n] = Map[Factor, Apart[FullSimplify[h[n]/h[n - 1]], \[Zeta]]];
    \[Alpha][n] =
     Map[Factor,
      Apart[
       FullSimplify[
        FullSimplify[-I/h[n]*
        Integrate[x*P[n, x]^2*w[x], {x, -1, 1}]] /. {BesselJ[2+\[Lambda], \[Zeta]] ->2 (\[Lambda] + 1)/\[Zeta] * BesselJ[\[Lambda] + 1,\[Zeta]] -
             BesselJ[\[Lambda], \[Zeta]]} /. {BesselJ[\[Lambda] + 1, \[Zeta]] -> (2 \[Lambda]/\[Zeta]) BesselJ[\[Lambda], \[Zeta]]}], \[Zeta]]]];
   {n, \[Alpha][n], \[Beta][n]}, {n, 1, 2}];

(*output the values of \[Alpha][n] and \[Beta][n] for n=0,1,2 in tabular form*)
TableForm[Join[{{0, \[Alpha][0], \[Beta][0]}}, results],
 TableHeadings -> {None, {"n", "\[Alpha][n]", "\[Beta][n]"}}]
 \end{lstlisting}

\section{Explicit forms of $\{X_i, Y_i, W_i,G_i\}$ from Tables \ref{Table-1}-\ref{tab2:rec-coeffs}.}\label{symbols}
(1) Explicit forms of $X_{i}(\zeta^2)$ from Table \ref{Table-1} for $i=2,3,4,5,6,8,9$.
\begin{align*}
X_{2}:=&\zeta ^4-\zeta ^2 \lambda  (2 \lambda +5)+4 \lambda ^2,\\
2X_{3}:=&2 \zeta ^6-\zeta ^4 \left(2 \lambda ^2+11 \lambda +3\right)+2 \zeta ^2 \lambda  \left(2 \lambda ^2+17 \lambda +12\right)-8 \lambda ^2 \left(2 \lambda ^3+7 \lambda ^2+11 \lambda +6\right),\\
4X_{4}:=&4 \zeta ^8-\zeta ^6 \left[4 \lambda  (\lambda +9)+21\right]+12 \zeta ^4 \lambda  [4 \lambda  (\lambda +5)+15]-16 \zeta ^2 \lambda ^2 (\lambda +1) [4 \lambda  (\lambda +4)+27]+192 \lambda ^3 (\lambda +1),\\
8X_{5}:=&8 \zeta ^{10}-5 \zeta ^8 \left(16 \lambda ^2+54 \lambda +33\right)+2 \zeta ^6 \left(68 \lambda ^4+548 \lambda ^3+1233 \lambda ^2+774 \lambda +36\right)\\
&-16 \zeta ^4 \lambda  (\lambda +1) \left(4 \lambda ^4+42 \lambda ^3+188 \lambda ^2+279 \lambda +54\right)+192 \zeta ^2 \lambda ^2 (\lambda +1) (\lambda +2) (14 \lambda +9)\\
&-384 \lambda ^3 (\lambda +1)^2 \left(2 \lambda ^3+7 \lambda ^2+12 \lambda +12\right),\\
8X_{6}:=&8 \zeta ^{12}-24 \zeta ^{10} \left(2 \lambda ^2+9 \lambda +6\right)+3 \zeta ^8 \left(24 \lambda ^4+228 \lambda ^3+674 \lambda ^2+527 \lambda +66\right)\\
&-4 \zeta ^6 \lambda  (\lambda +1) \left(8 \lambda ^4+100 \lambda ^3+590 \lambda ^2+1431 \lambda +612\right)+144 \zeta ^4 \lambda ^2 (\lambda +1) (2 \lambda +9) \left(2 \lambda ^2+11 \lambda +8\right)\\
&-192 \zeta ^2 \lambda ^3 (\lambda +1)^2 (2 \lambda +7) \left(2 \lambda ^2+5 \lambda +12\right)+4608 \lambda ^4 (\lambda +1)^2,\\
16X_{8}:=&16 \zeta ^{16}-4 \zeta ^{14} \left(12 \lambda ^2+88 \lambda +81\right)+\zeta ^{12} \left(48 \lambda ^4+832 \lambda ^3+4496 \lambda ^2+6736 \lambda +2889\right)\\
&-\zeta ^{10} \left(16 \lambda ^6+1632 \lambda ^5+17128 \lambda ^4+72456 \lambda ^3+128917 \lambda ^2+94824 \lambda +23076\right)\\
&+8 \zeta ^8 (\lambda +1) \left(400 \lambda ^6+5568 \lambda ^5+29576 \lambda ^4+76800 \lambda ^3+91173 \lambda ^2+37476 \lambda +540\right)\\
&-64 \zeta ^6 \lambda  (\lambda +1) \left(48 \lambda ^7+824 \lambda ^6+5924 \lambda ^5+22426 \lambda ^4+47488 \lambda ^3+51882 \lambda ^2+22527 \lambda +1080\right)\\
&+64 \zeta ^4 \lambda ^2 (\lambda +1)^2 \left(16 \lambda ^7+288 \lambda ^6+2344 \lambda ^5+11520 \lambda ^4+34117 \lambda ^3+57726 \lambda ^2+44892 \lambda +6480\right)\\
&+3072 \zeta ^2 \lambda ^3 (\lambda +1)^2 \left(4 \lambda ^5-16 \lambda ^4-253 \lambda ^3-842 \lambda ^2-1044 \lambda -360\right)\\
&+6144 \lambda ^4 (\lambda +1)^3 \left(4 \lambda ^5+32 \lambda ^4+107 \lambda ^3+202 \lambda ^2+258 \lambda +180\right),\\
16X_{9}:=&16 \zeta ^{18}-12 \zeta ^{16} \left(4 \lambda ^2+36 \lambda +41\right)+3 \zeta ^{14} \left(16 \lambda ^4+480 \lambda ^3+3000 \lambda ^2+5272 \lambda +2601\right)\\
&-\zeta ^{12} \left(16 \lambda ^6+3504 \lambda ^5+40344 \lambda ^4+176296 \lambda ^3+324993 \lambda ^2+246267 \lambda +61254\right)\\
&+12 \zeta ^{10} (\lambda +1) \left(496 \lambda ^6+7456 \lambda ^5+43448 \lambda ^4+122600 \lambda ^3+157107 \lambda ^2+70500 \lambda +2880\right)\\
&-24 \zeta ^8 \lambda  (\lambda +1) \left(208 \lambda ^7+3872 \lambda ^6+30552 \lambda ^5+129520 \lambda ^4+310013 \lambda ^3+384306 \lambda ^2+198840 \lambda +23400\right)\\
&+96 \zeta ^6 \lambda ^2 (\lambda +1)^2 \left(16 \lambda ^7+320 \lambda ^6+2904 \lambda ^5+16144 \lambda ^4+57113 \lambda ^3+121224 \lambda ^2+123780 \lambda +36000\right)\\
&-4608 \zeta ^4 \lambda ^3 (\lambda +1)^2 \left(8 \lambda ^5+180 \lambda ^4+1214 \lambda ^3+3687 \lambda ^2+4880 \lambda +2100\right)\\
&+9216 \zeta ^2 \lambda ^4 (\lambda +1)^3 \left(8 \lambda ^5+84 \lambda ^4+374 \lambda ^3+939 \lambda ^2+1526 \lambda +1200\right)-1105920 \lambda ^5 (\lambda +1)^3 (\lambda +2).
\end{align*}
(2) Explicit forms of $Y_{i}(\zeta^2)$ from Table \ref{Table-1-0} for $i=2,3,4,5,6,8,9,11,12,15,16$.
\begin{align*}
4Y_{2}:=&4 \zeta ^4-72 \zeta ^2+99,\\
16Y_{3}:=&16 \zeta ^6-300 \zeta ^4+2187 \zeta ^2-17622,\\
16Y_{4}:=&16 \zeta ^8-492 \zeta ^6+7803 \zeta ^4-61254 \zeta ^2+34560,\\
96Y_{5}:=&96 \zeta ^{10}-11944 \zeta ^8+291822 \zeta ^6-2328885 \zeta ^4+6414408 \zeta ^2-26226720,\\
64Y_{6}:=&64 \zeta ^{12}-6144 \zeta ^{10}+153144 \zeta ^8-1548729 \zeta ^6+8883864 \zeta ^4-40214880 \zeta ^2+11404800,\\
768Y_{8}:=&768 \zeta ^{16}-62784 \zeta ^{14}+2295136 \zeta ^{12}-83342508 \zeta ^{10}+2029268835 \zeta ^8-23810490144 \zeta ^6\\
&+113854569840 \zeta ^4-166228536960 \zeta ^2+405791078400,\\
256Y_{9}:=&256 \zeta ^{18}-26496 \zeta ^{16}+1540080 \zeta ^{14}-68175864 \zeta ^{12}+1655096733 \zeta ^{10}-19800590616 \zeta ^8\\
&+111180042720 \zeta ^6-320921429760 \zeta ^4+944392089600 \zeta ^2-151538688000,\\
12288Y_{11}:=&12288\zeta^{22}-4841984 \zeta ^{20}+498115584 \zeta ^{18}-22487382864 \zeta ^{16}+544353295644 \zeta ^{14}\nonumber\\
&-8913114152811 \zeta ^{12}+129651131201400 \zeta ^{10}-1583714648570400 \zeta ^8+11730625253222400 \zeta ^6\nonumber\\
&-38412263198822400 \zeta ^4+34441015482777600 \zeta ^2-55575779180544000,\\
256Y_{12}:=&256\zeta ^{24}-76800 \zeta ^{22}+7503840 \zeta ^{20}-351447300 \zeta ^{18}+9849042720 \zeta ^{16}-211834514439 \zeta ^{14}\nonumber\\
&+3852242371395 \zeta ^{12}-50175327507750 \zeta ^{10}+381546799239600 \zeta ^8-1420808577840000 \zeta ^6\nonumber\\
&+2410899503270400 \zeta ^4-4958458219008000 \zeta ^2+487223009280000,\\
16384Y_{15}:=&16384\zeta ^{30}-3735552 \zeta ^{28}+399025152 \zeta ^{26}-37752509952 \zeta ^{24}+3191338360512 \zeta ^{22}\nonumber\\
&-180668317734288 \zeta ^{20}+6264700943957640 \zeta ^{18}-134725503695291055 \zeta ^{16}\nonumber\\
&+1865006354944081260 \zeta ^{14}-18155774036247088680 \zeta ^{12}+149392043600809982400 \zeta ^{10}\nonumber\\
&-1120981072745389824000 \zeta ^8+5857439245474692096000 \zeta ^6-14118903261833158656000 \zeta ^4\nonumber\\
&+8447643101325754368000 \zeta ^2-9665272382965678080000,\\
16384Y_{16}:=&16384 \zeta ^{32}-4300800 \zeta ^{30}+627102720 \zeta ^{28}-76294333440 \zeta ^{26}+6704550057600 \zeta ^{24}\nonumber\\
&-371834537851296 \zeta ^{22}+12845996840802960 \zeta ^{20}-285997297566804075 \zeta ^{18}\nonumber\\
&+4385348005754977200 \zeta ^{16}-52253512139882473200 \zeta ^{14}+541564860087827107200 \zeta ^{12}\nonumber\\
&-4545805285706008320000 \zeta ^{10}+24640777114228684800000 \zeta ^8-66600404891865169920000 \zeta ^6\nonumber\\
&+72962121079040901120000 \zeta ^4-110368487460632002560000 \zeta ^2+7033911643629158400000.
\end{align*}
(3) Explicit forms of $W_i(\eta^2)$ from Table \ref{tab1:rec-coeffs} for $i=2,3,4,5,6,8,9,11,12,15,16$.
\begin{align*}
2W_2:=	&2\eta^4-15\eta^2+4,\\
2W_3:=	 &2\eta^6-8\eta^4+11\eta^2-168,\\
4W_4:=	 &4\eta^8-37\eta^6+295\eta^4-1248\eta^2+144,\\
8W_{5}:=&8\eta^{10}-773\eta^8+6761\eta^6-12522\eta^4-3072\eta^2-58176,\\
8W_{6}:=&8\eta^{12}-440\eta^{10}+3905\eta^8-11727\eta^6+32340\eta^4-110880\eta^2+6912,\\
16W_{8}:=&16\eta^{16}-436\eta^{14}+4129\eta^{12}-242046\eta^{10}+3296421\eta^8-14682816\eta^6+16861824\eta^4\\
&+8259840\eta^2+26542080,\\
16W_{9}:=&16\eta^{18}-636 \eta^{16}+21099\eta^{14}-723251\eta^{12}+8285121\eta^{10}-36265401\eta^8+56679264\eta^6\\
&-61715520\eta^4+216483840\eta^2-8294400,\\
16W_{11}:=&16\eta^{22}-5236\eta^{20}+268879\eta^{18}-4500486 \eta^{16}+26032641\eta^{14}-238158426\eta^{12}\\
&+4432046004\eta^{10}-34557365880\eta^8+105003993600\eta^6-85274726400\eta^4\\
&-41286205440\eta^2-56335564800,\\
64W_{12}:=&64\eta^{24}-12576\eta^{22}+603648\eta^{20}-10953887\eta^{18}+116382591\eta^{16}-2014897041\eta^{14}\\
&+27370556433\eta^{12}-179369958168\eta^{10}+520642779840\eta^8-515759892480\eta^6\\
&+190938746880\eta^4-932224204800\eta^2+23887872000,\\
256W_{15}:=&256\eta^{30}-24768\eta^{28}+984608\eta^{26}-99913860\eta^{24}+7272545505\eta^{22}-229238078161\eta^{20}\\
&+3336622422557\eta^{18}-22813872410709\eta^{16}+71457280019892\eta^{14}-369746460729216\eta^{12}\\
&+4741162121583360\eta^{10}-27852122634854400\eta^8+65627975905689600\eta^6\\
&-40565849142067200\eta^4-17279072442777600\eta^2-12917032353792000,\\
256W_{16}:=&256\eta^{32}-30592\eta^{30}+2610288 \eta^{28}-275214104\eta^{26}+16234800485 \eta^{24}-466147159164\eta^{22}\\
&+6682602532554 \eta^{20}-50829324276348\eta^{18}+308250455440113 \eta^{16}\\
&-3063456569159424\eta^{14}+26975907910160640 \eta^{12}-127620988169748480 \eta^{10}\\
&+279555488146022400\eta^8-196728971172249600\eta^6+7551166552473600\eta^4\\
&-185311791415296000 \eta^2+3371056496640000.
\end{align*}
(4) Explicit forms of $G_i(\eta^2)$ from Table \ref{tab2:rec-coeffs} for $i=2,3,4,5,6,8,9,11,12,15,16$.
\begin{align*}
G_{2}:=&\eta^4-13\eta^2+6,\\
G_{3}:=&\eta^6-7\eta^4+12\eta^2-360,\\
G_{4}:=&\eta^8-14\eta^6+165\eta^4-1224\eta^2+216,\\
G_{5}:=&\eta^{10}-133\eta^8+1929\eta^6-5940\eta^4-4320\eta^2-64800,\\	 G_{6}:=&\eta^{12}-77\eta^{10}+1101\eta^8-4932\eta^6+17361\eta^4-113400\eta^2+9720,\\	 2G_{8}:=&2\eta^{16}-78\eta^{14}+1089 \eta^{12}-70785\eta^{10}+1578285\eta^8-11659950 \eta^6+21402225 \eta^4\\
&+22963500\eta^2+76545000,\\
G_{9}:=&\eta^{18}-54 \eta^{16}+2286\eta^{14}-105966\eta^{12}+1927125\eta^{10}-13601520 \eta^8+31122225 \eta^6\\
&-30362850\eta^4+273375000\eta^2-13122000,\\
2G_{11}:=&2\eta^{22}-829\eta^{20}+58275\eta^{18}-1433430\eta^{16}+13530150\eta^{14}-117734175\eta^{12}\\
&+3083335875\eta^{10}-41419987875\eta^8+211076026875\eta^6-264998790000\eta^4\\
&-273265650000\eta^2-385786800000,\\
G_{12}:=&\eta^{24}-252\eta^{22}+16437\eta^{20}-425835\eta^{18}+6134265\eta^{16}-120820950\eta^{14}+2458605150\eta^{12}\\
&-26341594875\eta^{10}+124060399875\eta^8-176784778125\eta^6-8439086250\eta^4\\
&-699238575000\eta^2+20667150000,\\
G_{15}:=&\eta^{30}-126\eta^{28}+6654\eta^{26}-749439\eta^{24}+71087220\eta^{22}-3142068435\eta^{20}\\
&+67831410825\eta^{18}-727590460275\eta^{16}+3449814897375\eta^{14}-12254173249500\eta^{12}\\
&+266877060778125\eta^{10}-2987052554075625\eta^8+12107364149475000\eta^6\\
&-11246939027100000\eta^4-10207918728000000\eta^2-7655939046000000,\\
G_{16}:=&\eta^{32}-152\eta^{30}+15804\eta^{28}-2021700\eta^{26}+157758030\eta^{24}-6308511300\eta^{22}\\
&+132164631450\eta^{20}-1493816296500\eta^{18}+10920101842125\eta^{16}-117706861125000\eta^{14}\\
&+1646298202263750\eta^{12}-13233804005362500\eta^{10}+47682964163165625\eta^8\\
&-47494695181500000\eta^6-22188045568500000\eta^4-93572588340000000\eta^2\\
&+1822842630000000.
		\end{align*}


\section{When $\gamma=0$, Mathematica code to compute $\alpha_n$ and $\beta_n$ (for $n=0,1,2,3,4$) given in Table \ref{tab1:rec-coeffs}, with $\beta_0$ arbitrary.}\label{A4}
\begin{lstlisting}[language=Mathematica]
(*initialization*)
P[0, x_] = 1;
P[-1, x_] = 0;
w[x_] = x*(1 - x^2)^(-1/2)*Exp[I*\[Eta] x];
$Assumptions = BesselJ[0, \[Eta]] == 0 && \[Eta] > 0;
(**when \[Gamma]=1/2 and \[Eta]=m \[Pi], with m being an arbitrary nonzero positive integer, the above two lines of code should be replaced by
      w[x_] = x*Exp[I*\[Eta] x];
      $Assumptions = {m \[Element] Integers, m > 0};
      \[Eta] = m*\[Pi];**)

(*definitions of h[0] and \[Alpha][0]*)
h[0] = Integrate[w[x], {x, -1, 1}];
\[Alpha][0] = Simplify[-I/h[0]*Integrate[x*w[x], {x, -1, 1}] /.
              BesselJ[2, \[Eta]] -> (2/\[Eta]) BesselJ[1, \[Eta]]];
(**when \[Gamma]=1/2, skip the substitution in the previous line**)

(*calculations of \[Alpha][n] and \[Beta][n] for n=1,2,3,4*)
results =
  Table[If[n > 0,
    P[n, x_] = (x - I*\[Alpha][n - 1])*P[n - 1, x] - \[Beta][n - 1]*P[n - 2, x];
    h[n] = Simplify[ Integrate[P[n, x]^2*w[x], {x, -1, 1}] /.
       BesselJ[2, \[Eta]] -> (2/\[Eta]) BesselJ[1, \[Eta]]];
  (**when \[lambda]=1/2, skip the substitution in the previous line**)
    \[Beta][n] = h[n]/h[n - 1];
    \[Alpha][n] =
     Map[Factor,
      Apart[
       Simplify[
         Simplify[-I/h[n]*Integrate[x*P[n, x]^2*w[x], {x, -1, 1}] /.
            BesselJ[2, \[Eta]] -> (2/\[Eta]) BesselJ[1, \[Eta]]] /.
          BesselJ[1, x] -> j] /. j -> BesselJ[1, x],\[Eta]]]];
   {n, \[Alpha][n], \[Beta][n]}, {n, 1, 4}];
   (**when \[Gamma]=1/2, skip the substitutions in the previous line**)

(**when \[Gamma]=1/2 and \[zeta]=m \[Pi], the following two lines of code need to be evaluated.
  Clear[\[Eta]];
  resultsFixed = results //. {m^k_*\[Pi]^k_ -> \[Eta]^k, m*\[Pi] -> \[Eta]};**)

(*output the values of \[Alpha][n] and \[Beta][n] for n=0,1,2,3,4 in tabular form*)
TableForm[Join[{{0, \[Alpha][0], \[Beta][0]}}, results],TableHeadings -> {None, {"n", "\[Alpha][n]", "\[Beta][n]"}}]
(**when \[Gamma]=1/2, replace "results" with "resultsfixed" in the line above**)
\end{lstlisting}

\section{When $\gamma=0$, an alternative derivation of the initial values $\{\alpha_0,\alpha_1,\alpha_2,\alpha_3\}$ and $\{\beta_1,\beta_2,\beta_3,\beta_4\}$ given in Table \ref{tab1:rec-coeffs}.}\label{J0D-int}
We need the following lemma.
\begin{lemma}
Supposing $\eta$ is a positive zero of the Bessel function $J_0$, we have the following six integral identities:
\begin{align*}
\frac{1}{\pi}\int_{-1}^{1}x(1-x^2)^{-\frac{1}{2}}\exp(\mathit{i}\eta x)\, dx=&i J_1(\eta),\\
\frac{1}{\pi}\int_{-1}^{1}x^{2}(1-x^2)^{-\frac{1}{2}}\exp(\mathit{i}\eta x)\, dx=&-\frac{J_1(\eta)}{\eta},\\
\frac{1}{\pi}\int_{-1}^{1}x^{3}(1-x^2)^{-\frac{1}{2}}\exp(\mathit{i}\eta x)\, dx=&i\left(-\frac{2}{\eta^2} +1\right) J_1(\eta),\\
\frac{1}{\pi}\int_{-1}^{1}x^{4}(1-x^2)^{-\frac{1}{2}}\exp(\mathit{i}\eta x)\, dx=&2\left(\frac{3}{\eta^3}-\frac{1}{\eta}\right)J_1(\eta),\\
\frac{1}{\pi}\int_{-1}^{1}x^{5}(1-x^2)^{-\frac{1}{2}}\exp(\mathit{i}\eta x)\, dx=&i\left(\frac{24}{\eta^4}-\frac{7}{\eta^2}+1\right)J_1(\eta),\\
\frac{1}{\pi}\int_{-1}^{1}x^{6}(1-x^2)^{-\frac{1}{2}}\exp(\mathit{i}\eta x)\, dx=&-3\left(\frac{40}{\eta^5}-\frac{11}{\eta^3}+\frac{1}{\eta}\right)J_1(\eta),\\
\frac{1}{\pi}\int_{-1}^{1}x^{7}(1-x^2)^{-\frac{1}{2}}\exp(\mathit{i}\eta x)\, dx=&i \left(-\frac{720}{\eta^6}+\frac{192}{\eta^4}-\frac{15}{\eta^2}+1\right)J_1(\eta),\\
\frac{1}{\pi}\int_{-1}^{1}x^{8}(1-x^2)^{-\frac{1}{2}}\exp(\mathit{i}\eta x)\, dx=&4\left(\frac{1260}{\eta ^7}-\frac{330}{\eta ^5}+\frac{24}{\eta
   ^3}-\frac{1}{\eta }\right) J_1(\eta ),\\
\frac{1}{\pi}\int_{-1}^{1}x^{9}(1-x^2)^{-\frac{1}{2}}\exp(\mathit{i}\eta x)\, dx=&i\left(\frac{40320}{\eta ^8}-\frac{10440}{\eta ^6}+\frac{729}{\eta
   ^4}-\frac{26}{\eta ^2}+1\right) J_1(\eta ).
%   \frac{1}{\pi}\int_{-1}^{1}x^{10}(1-x^2)^{-\frac{1}{2}}\exp(\mathit{i}\eta x)\, dx=&-5\left(\frac{72576}{\eta ^9}-\frac{18648}{\eta ^7}+\frac{1269}{\eta
%   ^5}-\frac{42}{\eta ^3}+\frac{1}{\eta }\right) J_1(\eta).
\end{align*}
\end{lemma}

\begin{proof}
From \cite[eq. 10, p. 912]{GR}, we know that
\begin{align}\label{bsint}
\frac{1}{\pi}\int_{-1}^1(1-x^2)^{-\frac{1}{2}}\exp(i\eta x)dx=J_0(\eta).
\end{align}
Differentiating it $m$ times with respect to $\eta$ yields
\begin{align}\label{DJ0m}
\frac{i^m}{\pi}\int_{-1}^{1}x^{m}(1-x^2)^{-\frac{1}{2}}\exp(\mathit{i}\eta x)\, dx=\frac{d^m}{d\eta^m}J_0(\eta)=:J_0^{(m)}(\eta).
\end{align}
Recall the following properties of the Bessel function \cite[(5.3.7)]{Lebedev}:
\begin{align*}
\frac{d}{dz}(z^{\nu}J_{\nu}(z))=&z^{\nu}J_{\nu-1}(z),&\frac{d}{dz}(z^{-\nu}J_{\nu}(z))=-z^{-\nu}J_{\nu+1}(z).
\end{align*}
With $\nu=1$ in the first identity and $\nu=0$ in the second one, we have
\begin{align}
\frac{d}{dz}(J_{1}(z))=&-\frac{J_1(z)}{z}+J_{0}(z),\label{DJ1}\\
\frac{d}{dz}(J_{0}(z))=&-J_{1}(z).\label{DJ0}
\end{align}
Setting $z=\eta$ in \eqref{DJ0} yields
\[J_0'(\eta)=-J_1(\eta).\]
 Differentiating \eqref{DJ0} with respect to $z$, in light of \eqref{DJ1}, we get
\begin{align}\label{DJ0-2}
\frac{d^2}{dz^2}(J_{0}(z))=\frac{J_1(z)}{z}-J_{0}(z).
\end{align}
In view of $J_0(\eta)=0$, it follows that
\[J_{0}''(\eta)=\frac{J_1(\eta)}{\eta}.\]
Taking the derivative of \eqref{DJ0-2} with respect to $z$, on account of \eqref{DJ1}-\eqref{DJ0}, we produce
\begin{align}\label{DJ0-3}
\frac{d^3}{dz^3}(J_{0}(z))=\left(-\frac{2}{z^2}+1\right)J_1(z)+\frac{J_0(z)}{z},
\end{align}
so that
\begin{align*}
J_{0}'''(\eta)=\left(-\frac{2}{\eta^2}+1\right)J_1(\eta).
\end{align*}
Following the same reasoning as in the determination of $J_0''(\eta)$ and $J_0'''(\eta)$, we obtain
\begin{align*}
J_0^{(4)}(\eta)=&\left(\frac{6}{\eta^3}-\frac{2}{\eta}\right)J_1(\eta),\\
J_0^{(5)}(\eta)=&\left(-\frac{24}{\eta^4}+\frac{7}{\eta^2}-1\right)J_1(\eta),\\
J_0^{(6)}(\eta)=&\left(\frac{120}{\eta^5}-\frac{33}{\eta^3}+\frac{3}{\eta}\right)J_1(\eta),\\
J_0^{(7)}(\eta)=&\left(\frac{-720}{\eta^6}+\frac{192}{\eta^4}-\frac{15}{\eta^2}+1\right)J_1(\eta),\\
J_0^{(8)}(\eta)=&\left(\frac{5040}{\eta ^7}-\frac{1320}{\eta ^5}+\frac{96}{\eta
   ^3}-\frac{4}{\eta }\right) J_1(\eta ),\\
J_0^{(9)}(\eta)=&\left(-\frac{40320}{\eta ^8}+\frac{10440}{\eta ^6}-\frac{729}{\eta
   ^4}+\frac{26}{\eta ^2}-1\right) J_1(\eta).
%   J_0^{(10)}(\eta)=&\left(\frac{362880}{\eta ^9}-\frac{93240}{\eta ^7}+\frac{6345}{\eta
%   ^5}-\frac{210}{\eta ^3}+\frac{5}{\eta }\right) J_1(\eta).
\end{align*}
Substituting $J_0^{(m)}(\eta)$ for $1\leq m\leq 9$ into \eqref{DJ0m} leads us to the desired integral formulas.
\end{proof}

Using \eqref{pf0-eq2}-\eqref{pf0-eq5}, with the aid of the integral identities listed in the above lemma, by a reasoning similar to the one given below Lemma \ref{w1L}, we obtain the values of $\alpha_n$  for $n=0,1,2,3$ and $\beta_n$ for $n=1,2,3,4$.

%\section{When $\lambda=1/2$, Mathematica code to compute $\alpha_n$ and $\beta_n$ for $n=0,1,2,3,4$, with $\beta_0=A$ given.}\label{A5}
%\begin{lstlisting}[language=Mathematica]
%(*initialization*)
%P[0, x_] = 1;
%P[-1, x_] = 0;
%\[Beta][0] = A;
%w[x_] = x*Exp[I*\[Zeta] x];
%$Assumptions = {m \[Element] Integers, m > 0};
%\[Zeta] = m*\[Pi];
%
%(*definitions of h[0] and \[Alpha][0]*)
%h[0] = Integrate[w[x], {x, -1, 1}];
%\[Alpha][0] = -I/h[0]*Integrate[x*w[x], {x, -1, 1}];
%
%(*calculations of \[Alpha][n] and \[Beta][n] for n=1,2,3,4*)
%results =
%Table[
%If[n > 0, P[n,x_]=(x - I*\[Alpha][n - 1])*P[n - 1, x] - \[Beta][n - 1]*P[n - 2,x];
%   h[n] = Integrate[P[n, x]^2*w[x], {x, -1, 1}];
%   \[Beta][n] = h[n]/h[n - 1];
%   \[Alpha][n] =
%     Map[Factor, Apart[-I/h[n]*Integrate[x*P[n, x]^2*w[x], {x, -1, 1}]]]];
%{n, \[Alpha][n], \[Beta][n]}, {n, 0, 4}];
%
%(*substitution of \[Zeta] for m\[Pi] in the obtained expressions of
%\[Alpha][n] and \[Beta][n]*)
%Clear[\[Zeta]];
%resultsFixed = results //. {m^k_*\[Pi]^k_ -> \[Zeta]^k, m*\[Pi] -> \[Zeta]};
%
%(*output the values of \[Alpha][n] and \[Beta][n] for n=0,1,2,3,4 in tabular form.*)
%TableForm[resultsFixed, TableHeadings -> {None, {"n", "\[Alpha][n]", "\[Beta][n]"}}]
%
%\end{lstlisting}

\section{When $\gamma=1/2$ and $\eta=100\pi$, Mathematica code to compute $\alpha_n$ and $\beta_n$ for $0\leq n\leq30$, with $\beta_0$ arbitrary.}\label{A6}

\begin{lstlisting}[language=Mathematica]
Clear[\[Gamma], \[Alpha], \[Beta], \[Eta]];
\[Gamma] = 1/2;

(*the initial values of \[Alpha][n] for n=0,1,2,3*)
\[Alpha][0] = 2/\[Eta];
\[Alpha][1] = 4/\[Eta] - (4 \[Eta])/(-2 + \[Eta]^2);
\[Alpha][2] = 6/\[Eta] + (4 \[Eta])/(-2 + \[Eta]^2) - (2 \[Eta] (-24 +
              \[Eta]^2))/(6 - 13 \[Eta]^2 + \[Eta]^4);
\[Alpha][3] = 8/\[Eta] + (2 \[Eta] (-24 + \[Eta]^2))/(6 - 13 \[Eta]^2 + \[Eta]^4) - (12 \[Eta] (-360 + 12 \[Eta]^2 - 7 \[Eta]^4 + \[Eta]^6))/(216 - 1224 \[Eta]^2 + 165 \[Eta]^4 - 14 \[Eta]^6 + \[Eta]^8);

(*the initial values of \[Beta][n] for n=1,2,3,4, with \[Beta][0] arbitrary*)
\[Beta][1] = (\[Eta]^2 - 2)/\[Eta]^2;
\[Beta][2] = -((4 (6 - 13 \[Eta]^2 + \[Eta]^4))/(\[Eta]^2 (-2 + \[Eta]^2)^2));
\[Beta][3] = ((-2 + \[Eta]^2) (216 - 1224 \[Eta]^2 + 165 \[Eta]^4 - 14 \[Eta]^6 + \[Eta]^8))/(\[Eta]^2 (6 - 13 \[Eta]^2 + \[Eta]^4)^2);
\[Beta][4] = -(16 (6 - 13 \[Eta]^2 + \[Eta]^4) (9720 + \[Eta]^2 (-113400 + \[Eta]^2 (17361 - 4932 \[Eta]^2 + 1101 \[Eta]^4 - 77 \[Eta]^6 + \[Eta]^8))))/(\[Eta]^2 (216 - 1224 \[Eta]^2 + 165 \[Eta]^4 - 14 \[Eta]^6 + \[Eta]^8)^2);

(*taking \[Eta]=100*\[Pi]*)
\[Eta] = N[100 \[Pi], 100];
(**for \[Eta]=\[Pi],replacing this line by \[Eta]=N[\[Pi],100]**)

(*computing \[Alpha][i] and \[Beta][i] for i=4,5,6,7,by using the
above initial values and the two difference equations for the
recurrence coefficients*)
Do[\[Alpha][n + 1] =
   1/(\[Eta] \[Beta][n + 1]) (-3 \[Alpha][n]^2 - 2 \[Alpha][-1 + n] \[Alpha][n] - \[Alpha][-1 + n]^2 + 2 (\[Beta][n] - 1) + \[Beta][-1 + n] (\[Eta] (2\[Alpha][-1 + n] + \[Alpha][-2 + n]) - 2 n - 2 \[Gamma] + 3) + (\[Alpha][-1 + n] -\[Alpha][           n]) ((\[Alpha][-2 + n] - \[Alpha][n]) (\[Eta] (\[Alpha][n] + \[Alpha][-1 + n] + \[Alpha][-2 + n]) - 2 n - 2 \[Gamma]) + \[Eta] (\[Beta][n] - \[Beta][-1 + n])) + (\[Alpha][-1 + n] - \[Alpha][n])/(\[Alpha][-2 + n] - \[Alpha][-1 + n]) (\[Alpha][-1 + n]^2 + 5 \[Alpha][-2 + n]^2 + 2 + \[Beta][n] (\[Eta] (\[Alpha][n] + 2 \[Alpha][-1 + n]) - 2 n - 2 \[Gamma] - 1) -  2 \[Beta][-1 + n] + \[Beta][-2 + n] (2 n + 2 \[Gamma] - 5 - \[Eta] (2 \[Alpha][-2 + n] + \[Alpha][-3 + n])))) - 2 \[Alpha][n] + (2 n + 2 \[Gamma] + 3)/\[Eta];

  \[Beta][n + 2] = \[Alpha][n + 1]^2 + \[Alpha][n]*\[Alpha][n + 1] + \[Alpha][n]^2 - \[Alpha][-1 + n]^2 - \[Beta][n + 1] + \[Beta][n] + \[Beta][-1 + n] - 2 (n + 2 + \[Gamma]) \[Alpha][n + 1]/\[Eta] - (2 n + 2 \[Gamma] + 5) \[Alpha][n]/\[Eta] + 2 (n + \[Gamma] - 1) \[Alpha][-1 + n]/\[Eta] + \[Alpha][n]/(\[Eta]*\[Beta][n + 1]) (\[Alpha][n]^2 + 1 + \[Beta][n] (2 n + 2 \[Gamma] - 1 - \[Eta] (\[Alpha][n] + \[Alpha][-1 + n]))) + (\[Beta][n + 1] - \[Beta][n])/(\[Eta]*\[Beta][n + 1] (\[Beta][-1 + n] - \[Beta][n])) (\[Beta][   n] (\[Eta] (\[Alpha][n]^2 + \[Alpha][n]*\[Alpha][-1 + n] - \[Beta][n + 1]) -          2 (n + 1 + \[Gamma]) \[Alpha][n] - 3 \[Alpha][-1 + n]) - \[Beta][-1 + n] (\[Eta] (\[Alpha][-2 + n]^2 + \[Alpha][-2 + n]*\[Alpha][-1 + n] - \[Beta][-2 + n]) - 2 (n + \[Gamma] - 2) \[Alpha][-2 + n] + 3 \[Alpha][-1 + n]) + \[Alpha][-1 + n]^3 + \[Alpha][-1 + n]);, {n, 3, 29}];

(*Format number with 18 significant digits;using 18 decimal places
for|x|>1*)
formatNumber[x_] :=
Module[{},If[Abs[x] > 1, NumberForm[x, {Infinity, 18}], NumberForm[x, 18]]];
(**When \[Eta]=\[Pi],18 is replaced by 16**)

(*Generating formatted table for \[Alpha][n] and \[Beta][n],0 \[LessEqual] n\[LessEqual]30; \[Beta][0] keeping unformatted,others using formatNumber*)
Grid[Prepend[Table[{n, formatNumber[FullSimplify[\[Alpha][n]]],
    If[n == 0, \[Beta][n], formatNumber[FullSimplify[\[Beta][n]]]]}, {n, 0, 30}], {"n", "\[Alpha][n]", "\[Beta][n]"}], Frame -> All]
\end{lstlisting}

\end{appendices}}

\end{document}

%------------------------------------------------------------------------------
% End of journal.tex
%------------------------------------------------------------------------------
